% concatenated from main.tex, sections/introduction.tex, sections/problemsetup.tex, sections/theory.tex, sections/experiments.tex, sections/conclusion.tex, appendices/cs4ml.tex, appendices/furtherexperimentdetails.tex, appendices/additionalexperiment.tex, appendices/cfgablation.tex, appendices/proofs.tex, appendices/recoveredimagecomps.tex
\documentclass{article}

% if you need to pass options to natbib, use, e.g.:
%     \PassOptionsToPackage{numbers, compress}{natbib}
% before loading neurips_2026

% The authors should use one of these tracks.
% Before accepting by the NeurIPS conference, select one of the options below.
% 0. "default" for submission
\PassOptionsToPackage{numbers,compress}{natbib}
\usepackage[preprint]{neurips_2026}
% the "default" option is equal to the "main" option, which is used for the Main Track with double-blind reviewing.
% 1. "main" option is used for the Main Track
%  \usepackage[main]{neurips_2026}
% 2. "position" option is used for the Position Paper Track
%  \usepackage[position]{neurips_2026}
% 3. "eandd" option is used for the Evaluations & Datasets Track
% \usepackage[eandd]{neurips_2026}
% if you want to opt-in for a de-anomymized submission (not recommended, but OK):
%\usepackage[eandd, nonanonymous]{neurips_2026}
% 4. "creativeai" option is used for the Creative AI Track
%  \usepackage[creativeai]{neurips_2026}
% 5. "sglblindworkshop" option is used for the Workshop with single-blind reviewing
% \usepackage[sglblindworkshop]{neurips_2026}
% 6. "dblblindworkshop" option is used for the Workshop with double-blind reviewing
%  \usepackage[dblblindworkshop]{neurips_2026}

% After being accepted, the authors should add "final" behind the track to compile a camera-ready version.
% 1. Main Track
% \usepackage[main, final]{neurips_2026}
% 2. Position Paper Track
%  \usepackage[position, final]{neurips_2026}
% 3. Evaluations & Datasets Track
% \usepackage[eandd, final]{neurips_2026}
% 4. Creative AI Track
%  \usepackage[creativeai, final]{neurips_2026}
% 5. Workshop with single-blind reviewing
%  \usepackage[sglblindworkshop, final]{neurips_2026}
% 6. Workshop with double-blind reviewing
%  \usepackage[dblblindworkshop, final]{neurips_2026}
% Note. For the workshop paper template, both \title{} and \workshoptitle{} are required, with the former indicating the paper title shown in the title and the latter indicating the workshop title displayed in the footnote.
% For workshops (5., 6.), the authors should add the name of the workshop, "\workshoptitle" command is used to set the workshop title.
% \workshoptitle{WORKSHOP TITLE}

% "preprint" option is used for arXiv or other preprint submissions
% \usepackage[preprint]{neurips_2026}

% to avoid loading the natbib package, add option nonatbib:
%    \usepackage[nonatbib]{neurips_2026}

\usepackage[utf8]{inputenc} % allow utf-8 input
\usepackage[T1]{fontenc}    % use 8-bit T1 fonts
\usepackage{hyperref}       % hyperlinks
\usepackage{url}            % simple URL typesetting
\usepackage{booktabs}       % professional-quality tables
\usepackage{amsfonts}       % blackboard math symbols
\usepackage{nicefrac}       % compact symbols for 1/2, etc.
\usepackage{microtype}      % microtypography
\usepackage{xcolor}         % colors
\usepackage{amsmath}
\usepackage{amsthm}
\usepackage{amssymb}
\usepackage{algorithm}
\usepackage{algorithmic}
\usepackage{graphicx}
\usepackage{subcaption}

\theoremstyle{plain}
\newtheorem{theorem}{Theorem}[section]
\newtheorem{lemma}[theorem]{Lemma}
\newtheorem{corollary}[theorem]{Corollary}
\newtheorem{proposition}[theorem]{Proposition}
\newtheorem{claim}[theorem]{Claim}

% definition-style environments
\theoremstyle{definition}
\newtheorem{definition}[theorem]{Definition}
\newtheorem{assumption}[theorem]{Assumption}
\newtheorem{example}[theorem]{Example}
\newtheorem{exercise}[theorem]{Exercise}
\newtheorem{question}[theorem]{Question}

% remark-style environments
\theoremstyle{remark}
\newtheorem{remark}[theorem]{Remark}
\newtheorem{notation}[theorem]{Notation}

% Note. For the workshop paper template, both \title{} and \workshoptitle{} are required, with the former indicating the paper title shown in the title and the latter indicating the workshop title displayed in the footnote.
\title{Active Learning for Conditional Generative Compressed Sensing}

% The \author macro works with any number of authors. There are two commands
% used to separate the names and addresses of multiple authors: \And and \AND.
%
% Using \And between authors leaves it to LaTeX to determine where to break the
% lines. Using \AND forces a line break at that point. So, if LaTeX puts 3 of 4
% authors names on the first line, and the last on the second line, try using
% \AND instead of \And before the third author name.

\author{%
  Alexander DeLise \\ %\thanks{Use footnote for providing further information about author (webpage, alternative address)---\emph{not} for acknowledging funding agencies.} \\
  Department of Scientific Computing \\
  Department of Mathematics \\
  Florida State University \\
  \texttt{ard22l@fsu.edu} \\
  % examples of more authors
  \And
  Nick Dexter \\
  Department of Scientific Computing \\
  Florida State University \\
  \texttt{nick.dexter@fsu.edu}
  % Affiliation \\
  % Address \\
  % \texttt{email} \\
  % \AND
  % Coauthor \\
  % Affiliation \\
  % Address \\
  % \texttt{email} \\
  % \And
  % Coauthor \\
  % Affiliation \\
  % Address \\
  % \texttt{email} \\
  % \And
  % Coauthor \\
  % Affiliation \\
  % Address \\
  % \texttt{email} \\
}

\newcommand*{\snorm}[2]{\left\vert\!\left\vert\!\left\vert #1 \right\vert\!\right\vert\!\right\vert_{#2}}

\begin{document}

\maketitle

\begin{abstract}
Generative compressed sensing uses the range of a pretrained generator as a nonlinear model for recovering structured signals from limited measurements. We study a conditional version of this problem for image recovery from subsampled Fourier measurements using prompt-conditioned generative models. Our framework separates two roles of conditioning: the prompt used to design the sampling distribution and the prompt used to define the recovery model. For ReLU and Lipschitz conditional generators, we prove stable recovery bounds showing that prompt-matched Christoffel sampling retains the same Christoffel complexity constant as existing near-optimal generative compressed sensing theory, while prompt mismatch incurs an explicit compatibility penalty. Experiments with Stable Diffusion show that prompts meaningfully reshape Christoffel sampling distributions and influence image recovery. Overall, our results suggest that prompts should be treated as design variables with distinct effects on sensing, approximation, and recovery.
\end{abstract}



\section{Introduction}
Compressed sensing~\cite{donoho2006compressed} is a popular signal processing technique to recover a signal $\mathbf{f}^* \in \mathbb{C}^n$ from an undersampled observation $\mathbf{y} \in \mathbb{C}^m$, with $m \ll n$, obtained via some forward process
\begin{equation}
    \mathbf{y} = \mathbf{A} \mathbf{f}^* + \mathbf{e}, \label{eq:CS}
\end{equation}
with numerous applications throughout science and engineering~\cite{chen2008prior,lustig2007sparse,duarte2008single,yuan2021snapshot,herrmann2012fighting}.
Here $\mathbf{A} \in \mathbb{C}^{m \times n}$ is the sampling matrix (with possibly random entries) and $\mathbf{e} \in \mathbb{C}^m$ is a vector of additive noise. 

The inverse problem in Equation~\ref{eq:CS} is underdetermined. Thus one typically assumes some application-driven structure, usually sparsity of $\mathbf{f}^*$ in some domain, to usher a unique solution~\cite{candes2006robust, candes2006stable, tibshirani1996regression, hastie2015statistical, pati1993orthogonal, tropp2007signal, chen2001atomic, adcock2021compressive}, provided the sampling matrix $\mathbf{A}$ satisfies some condition like the Restricted Isometry Property~\cite{donoho2006compressed, candes2006stable}. Generative compressed sensing, on the other hand, replaces the sparsity assumption by instead requiring that $\mathbf{f}^*$ lies close to the range of a generative model $G : \mathbb{R}^k \to \mathbb{R}^n$, as first introduced in~\cite{bora2017compressed}. In this regime, $\mathbf{f}^*$ is recoverable with high probability when the sampling matrix $\mathbf{A}$ with Gaussian entries satisfies the so-called Set-Restricted Eigenvalue Condition. Several subsequent studies investigate upper and lower error bounds for signal recovery in generative compressed sensing~\cite{kamath2019lower, kamath2020power, liu2020information} as well as recovery guarantees under the presence of outlier data~\cite{jalal2020robust} or with nonlinear sampling operators~\cite{chen2023unified}.

While these results establish a powerful alternative to sparsity-based recovery, they rely on assumptions about the measurement process that are often not met in practice. First, Gaussian measurements can be unrealistic: in many practical applications, sensing is dictated by hardware constraints and corresponds to structured, often subsampled unitary transforms, e.g., the discrete Fourier transform in MRI~\cite{lustig2007sparse, lustig2008compressed}. Second, even after restricting to such physically realizable operators, the choice of sampling pattern remains critical. Common strategies such as Gaussian or uniform subsampling are uninformed, as they fail to exploit the structure of the signal class being measured and may lead to suboptimal performance~\cite{adcock2017breaking, krahmer2015compressive, poon2017structure}.


Several studies propose sampling strategies for generative compressed sensing that exploit the geometry of the underlying generative model, a process called \emph{active learning}. For example,~\cite{adcock2023cs4ml} develops the Christoffel Sampling for Machine Learning (CS4ML) framework based on the \emph{generalized Christoffel function}, while~\cite{berk2023modeladapted, berk2022fourier, plan2025denoising} use \emph{local coherences} to design sampling distributions that place more mass on measurement coordinates that are most informative for distinguishing signals in the model class. Such active learning approaches relate closely to the well-studied \emph{leverage score sampling}, with various applications in active learning, standard regression, and imaging~\cite{chen2016statistical,derezinski2018leveraged,ordozgoiti2022generalized,erdelyi2020fourier,avron2019universal}, as well optimal sampling strategies for compressed sensing in general~\cite{adcock2021compressive, krahmer2015compressive, poon2017structure, candes2007compressive, krahmer2013stable, puy2011variable}. 

Despite these advances, existing work in generative compressed sensing has largely relied on unconditional generative priors, thereby omitting auxiliary information that may be available at acquisition or inference time. In many realistic settings, such information can sharply restrict the set of plausible signals. For example, in medical imaging, modality, anatomy, or patient metadata may improve reconstruction~\cite{chung2025contextmri}; similarly, in natural image reconstruction, subject, scene category, acquisition protocol, or other metadata may narrow the relevant signal class and improve reconstruction quality~\cite{chung2024prompttune, kim2025regularization}. Prompted or otherwise conditional generators provide a natural mechanism for incorporating this information by steering the recovery class toward signals compatible with the available context, rather than requiring a single unconditional model to represent all possible signal types.

Conditioning can also improve sensing itself. If the relevant conditional signal class is narrower, then one may design sampling distributions supported on more informative measurement coordinates, leading to lower sample complexity for stable recovery. However, conditioning introduces interactions absent from the unconditional setting: the prompt used for sampling need not match the prompt used for recovery, and neither need match the class containing the true signal, when $\mathbf f^*$ lies in the range of the generator. Thus, overly strong or misaligned conditioning may bias the reconstruction toward the prior and away from the observed measurements. This motivates \emph{conditional generative compressed sensing} as an extension of the standard framework in which both sampling design and recovery guarantees depend explicitly on the interaction between conditional signal classes. Understanding how these interactions affect stable recovery and reconstruction is the central question of this paper.

\subsection{Contributions}
We adapt an active learning perspective via Christoffel sampling from~\cite{adcock2023cs4ml} together with generative compressed sensing arguments from~\cite{bora2017compressed, berk2023modeladapted, berk2022fourier} to study prompt-conditioned signal recovery under subsampled Fourier measurements. The resulting theory characterizes how interactions between the sampling, recovery, and true-signal prompts determine the sample complexity required for stable recovery and the resulting reconstruction error. Our contributions may be summarized as follows:
\begin{itemize}
    \item We formulate a prompt-conditioned Christoffel sampling framework for generative compressed sensing under subsampled Fourier measurements with three conditioning prompts: the true-signal prompt $c_*$, the recovery prompt $c_r$, and the sampling prompt $c_s$.

    \item We identify a single prompt compatibility factor $\Lambda(c_1,c_2,c_3)$ and show that evaluating it at $(c_r,c_r,c_s)$ controls sample complexity for stable reconstruction, while evaluating it at $(c_*,c_r,c_s)$ controls the residual error in recovery.

    \item We derive explicit sample complexity results under ReLU and Lipschitz assumptions on the conditional generator and combine these with agnostic signal recovery error bounds.

    \item We empirically test the theoretical framework through a suite of experiments demonstrating that prompt conditioning systematically influences both the induced sampling distributions and reconstruction performance.
\end{itemize}

\section{Active Learning Problem Setup}\label{sec: problem setup}

We follow the general Christoffel sampling for machine learning (CS4ML) framework established in~\cite{adcock2023cs4ml} for our active learning regime, however we specialize our results to subsampled Fourier measurements. For a full description of the general CS4ML framework, we refer the reader to Appendix~\ref{app:cs4ml}.

\subsection{Setup and Prompt-Conditioned Model Classes}
Let $\mathbb{X} = \mathbb{R}^n$, deemed the \emph{object space}, be equipped with Euclidean norm $\left\| \cdot \right\|_2$. Let $\mathbf{F} \in \mathbb{C}^{n \times n}$ denote the unitary discrete Fourier transform, let $D = \{1, \dots, n\}$, and let $\mathbf{P}_i$ be the row selector matrix for the $i$th Fourier coefficient. For a signal $\mathbf{f} \in \mathbb{X}$, define
\[
  L_i(\mathbf{f}) = \mathbf{P}_i \mathbf{F}\mathbf{f} \in \mathbb{C}, \qquad i \in D
\]
to be the \emph{sampling operator}, which extracts the $i$th Fourier coefficient of $\mathbf{f}$. We assume that the full family of sampling operators $\{L_i\}_{i \in D}$ is nondegenerate on $\mathbb{X}$, i.e.\ that there exist constants $0 < \alpha \le \beta < \infty$ such that
\[
  \alpha \left\|\mathbf{f}\right\|_2^2
  \le
  \sum_{i \in D} |L_i(\mathbf{f})|^2
  \le
  \beta \left\|\mathbf{f}\right\|_2^2,
  \qquad
  \mathbf{f} \in \mathbb{X}.
\]
In our setting, this holds with $\alpha = \beta = 1$ by unitarity of the discrete Fourier transform $\mathbf{F}$. Later, after introducing subsampling, we will require a sampling analogue of this property, deemed \emph{empirical nondegeneracy}, however we defer this notion to Section~\ref{sec: theory}. These properties ensure that the measurement map is injective and norm-preserving on $\mathbb{X}$, so that differences between signals are reflected in the measurements, thereby enabling stable and distinguishable recovery.

Next, let $G:\mathbb{R}^k\times\mathcal C\to\mathbb{X}$ be a conditional generative model which takes a latent vector $\mathbf{z} \in \mathbb{R}^k$ and conditioning $c \in \mathcal{C}$ as input, and outputs a signal $\mathbf{f} = G(\mathbf{z},c) \in \mathbb{X}$. Here, $\mathcal C$ denotes an abstract conditioning space, which may represent labels, metadata, measurements, or text embeddings, though in this work, we specialize to the case where conditioning is given by text prompts. In our subsequent analyses, we further assume the latent vector input $\mathbf{z}$ has bounded norm and let $B_2^k(R) = \left\{ \mathbf{z} \in \mathbb{R}^k \, : \, \left \| \mathbf{z} \right \|_2 \leq R \right\}$. Many generative models (e.g.\ diffusion-based~\cite{song2021scorebased,ho2020denoising,sohl2015deep}, GANs~\cite{goodfellow2014generative,radford2015unsupervised}, and VAEs~\cite{kingma2014auto}) assume the entries of $\mathbf{z}$ are Gaussian, therefore with these models, this assumption omits only an exponentially unlikely number of inputs in practical applications~\cite{bora2017compressed}. Lastly, define
\[
  \mathbb{F}_c = \left\{ G(\mathbf{z}, c) : \mathbf{z} \in B_2^k(R)\right\} \subseteq \mathbb{X},
\]
termed the \emph{approximation space}, to be the range of the generative model conditioned on a prompt $c$. 

In practice, we distinguish three different prompts throughout our analyses. First, we consider the signal prompt $c_*$, which manifests when our quantity of interest $\mathbf{f}^*$ is explicitly in the range of the generative model, i.e. $\mathbf{f}^* \in \mathbb{F}_{c_*}$. In general, $c_*$ may be unknown, or $\mathbf{f}^*$ might be out of distribution, in which case we also consider the recovery prompt $c_r$ which defines the set of candidate signals used to reconstruct $\mathbf{f}^*$ using $G$, that is, using signals in $\mathbb{F}_{c_r}$. Finally, we consider the sampling prompt $c_s$ used to construct our sampling distribution $\mu_{c_s}$, which we introduce in the following subsection. 

As we will see, a mismatch $c_* \neq c_r$ is a \emph{signal-recovery mismatch} and enters as an approximation error, whereas a mismatch $c_s \neq c_r$ is a \emph{sampling-recovery mismatch} and enters through the sample complexity required for stable signal recovery. Likewise, a mismatch $c_* \neq c_s$ is a \emph{signal-sampling mismatch}, meaning that the measurements are designed using prompt information that may not reflect the true type of signal being acquired, and in our analyses this manifests through how well the resulting sampling law captures the features of the target signal relevant for recovery.

\subsection{Christoffel Sampling and Prompt Compatibility}
A central component of active learning is the principled selection of measurements according to an optimized sampling strategy, which may be determined prior to data collection or updated adaptively as observations are acquired, in contrast to fixed or uniform sampling. In our setting, this corresponds to selecting indices $i \in D$ at which to query the sampling operators $L_i$. The goal is to allocate measurements so that they are maximally informative for distinguishing elements of the model class under consideration.

One such active sampling scheme is \emph{Christoffel sampling}, which harnesses the \emph{generalized Christoffel function}~\cite{adcock2023cs4ml, adcock2025optimal}. For any nonempty set $\mathbb S\subseteq \mathbb X$, define the Christoffel function of $\mathbb S$ in our subsampled Fourier measurement setting by
\[
  K(\mathbb S)(i)
  =
  \sup_{\mathbf{h} \in \mathbb S\setminus\left\{ \mathbf{0} \right\}}
  \frac{\left| L_i(\mathbf{h}) \right|^2}{\left\| \mathbf{h} \right\|_2^2}
  =
  \sup_{\mathbf{h} \in \mathbb S\setminus\left\{ \mathbf{0} \right\}}
  \frac{\left| \mathbf{P}_i \mathbf{F}\mathbf{h} \right|^2}{\left\| \mathbf{h} \right\|_2^2},
  \qquad i \in D,
\]
as well as the related quantity $\kappa(\mathbb S) = \sum_{i=1}^n K(\mathbb S)(i)$. In particular, for a prompt $c \in \mathcal{C}$ we apply this construction to the self-difference class $\mathbb{F}_c - \mathbb{F}_c = \left\{ \mathbf{f}_1 - \mathbf{f}_2 : \mathbf{f}_1, \mathbf{f}_2 \in \mathbb{F}_c \right\}$. We refer to elements in such sets as \emph{secants}. Since recovery is performed over a prompt-conditioned class $\mathbb{F}_{c_r}$, the notion of informativeness is inherently model-dependent and governed by how signals in $\mathbb{F}_{c_r}$ differ from one another. In particular, stable recovery requires that measurements capture directions along which elements of $\mathbb{F}_{c_r}$ can vary, i.e.\ the geometry of the self-difference class $\mathbb{F}_{c_r} - \mathbb{F}_{c_r}$. This is quantified by the Christoffel function $K(\mathbb{F}_c - \mathbb{F}_c)(i)$, which measures how much secants of $\mathbb{F}_c$ can concentrate in the $i$th Fourier coordinate. Larger values therefore correspond to coordinates that are more informative for distinguishing elements of the prompt-conditioned model class.

Under the standard active learning setup, we assume we have freedom to query $L_i$ for any $i \in D$, which allows us to design model-adapted sampling distributions using the generalized Christoffel function. The corresponding Christoffel sampling law is the probability distribution
\begin{equation}\label{eq: christoffel sampling law}
  \mu_c(i)
  =
  \frac{K(\mathbb{F}_c - \mathbb{F}_c)(i)}{\kappa(\mathbb{F}_c - \mathbb{F}_c)},
  \qquad i \in D.
\end{equation}
Here, $\mu_c$ allocates more probability mass to Fourier coordinates that can carry more normalized secant energy. To make our subsequent analyses more straightforward, we impose the following assumption, which follows from~\cite[Assumption 2.2]{adcock2023cs4ml}:

\begin{assumption}\label{assump: full support}
    This Christoffel sampling law $\mu_c$ has full-support, i.e. that $\mu_c(i) > 0$ for every $c \in \mathcal C$ and for every $i \in D$.
\end{assumption}


In the prompt-matched case, where the same prompt is used for both recovery and measurement, i.e. $c_r = c_s$, one would sample according to $\mu_{c_r}$. However, our framework allows the recovery prompt $c_r$ and the sampling prompt $c_s$ used to construct $\mu_{c_s}$ to differ. To quantify these prompt interactions with a single statistic, we define the following quantity:

\begin{definition}[Prompt Compatibility Factor]
  Let $c_1,c_2,c_3 \in \mathcal{C}$. Define the \emph{prompt compatibility factor} by
  \[
    \Lambda(c_1,c_2,c_3)
    =
    \max_{i \in D}
    \frac{K(\mathbb{F}_{c_1}-\mathbb{F}_{c_2})(i)}{\mu_{c_3}(i)}.
  \]
  \label{def: prompt compatibility factor}
\end{definition}

\begin{remark}[Reduction to Known Sampling Complexity Measures]
In the prompt-matched case where we recover over the same sampling prompt ($c_r = c_s$), we have $\Lambda(c_r,c_r,c_r)=\kappa(\mathbb F_{c_r}-\mathbb F_{c_r})$, reducing to the Christoffel sample-complexity constant from~\cite{adcock2023cs4ml}. More generally, recovering under a prompt different than that used to create the sampling distribution requires evaluating $\Lambda(c_r,c_r,c_s)$, which under our Fourier sampling model recovers the model-adapted coherence sample-complexity constant in~\cite{berk2023modeladapted, plan2025denoising}, which indeed is a special case of $\kappa$. What makes $\Lambda$ unique in the agnostic signal recovery setting is that it extends the usual self-difference analysis to cross-difference classes. When the target signal $\mathbf f^*$ is generated with a prompt $c_*$, the prompt compatibility factor $\Lambda$ not only controls sample complexity through $\mathbb F_{c_r}-\mathbb F_{c_r}$, but also quantifies the complexity of $\mathbb F_{c_*}-\mathbb F_{c_r}$. As we later show, this yields direct control of the approximation error in recovery.
\end{remark}

In the sampling-recovery mismatched case, $\Lambda(c_r,c_r,c_s)$ might blow up if Fourier frequencies needed for the recovery geometry under $c_r$ receive little positive mass under the sampling law from $c_s$. Of course, this immediately suggests that one should always set $c_r = c_s$ to circumvent this issue, though this is not always feasible in practice. Acquisition and reconstruction may occur at different times and with different levels of information. For example, at sensing time, one may only know coarse metadata or an initial prompt, and must therefore design measurements using $c_s$, whereas at reconstruction time one may have access to a more refined prompt and recover over $\mathbb{F}_{c_r}$. Moreover, one may wish to design a single sampling distribution that is effective across a family of possible recovery prompts, so that $c_s$ is chosen not to match one particular $c_r$ (or $c_*$), but rather as a compromise or universal design for several downstream recovery classes. Our theory quantifies how these mismatches change both the sample complexity needed to stabilize recovery and the resulting reconstruction error.

%%%%%%%%%%%%%%%%%%%%
%%%%%% Theory %%%%%%
%%%%%%%%%%%%%%%%%%%%
\section{Main Theoretical Results}\label{sec: theory}
Throughout this section, let $I_1,\dots,I_m$ be drawn independently from a sampling distribution $\mu_{c_s}$, write $\Omega= \{ I_1,\dots,I_m \}$, and let $\mathbf P_\Omega \in \mathbb R^{m\times n}$ denote the sampling matrix whose $j$th row is $\mathbf P_{I_j}$. Define the weighted subsampled Fourier operator
\[
  \mathbf A_\Omega
  =
  \frac{1}{\sqrt{m}} \mathbf W_\Omega^{1/2}\mathbf P_\Omega \mathbf F,
  \qquad
  \mathbf W_\Omega
  =
  \operatorname{diag}\left(
    \frac{1}{\mu_{c_s}(I_1)},
    \dots,
    \frac{1}{\mu_{c_s}(I_m)}
  \right),
\]
where $\mathbf F$ denotes the unitary discrete Fourier transform. We observe noisy weighted measurements
$
  \mathbf y = \mathbf A_\Omega \mathbf f_* + \mathbf e,
$
where $\mathbf e = \frac{1}{\sqrt{m}}\mathbf W_\Omega^{1/2}\mathbf u$ for some
$\mathbf{u}\in\mathbb C^m$ denoting unweighted noise.

Next, we introduce the two near-isometry principles that allow for stable signal recovery. 

\begin{definition}[S-REC~\cite{bora2017compressed}]
  For $\mathbb{S} \subseteq \mathbb X$, we say that $\mathbf{A}_\Omega$ satisfies the Set-Restricted Eigenvalue Condition (S-REC) with parameters $(\gamma, q)$, denoted $\mathrm{S\text{-}REC}\left(\mathbb S,\gamma,q\right)$, if
  \[
    \left\|\mathbf A_\Omega\left(\mathbf f_1-\mathbf f_2\right)\right\|_2
    \geq
    \gamma \left\|\mathbf f_1-\mathbf f_2\right\|_2
    -
    q,
    \qquad
    \mathbf f_1,\mathbf f_2\in\mathbb{S}.
  \]
\end{definition}

\begin{definition}[Empirical Nondegeneracy]\label{def:empirical-nondegeneracy}
  Let $0<\tau<1$. We say that $\mathbf A_\Omega$ is empirically nondegenerate on $\mathbb{F}_{c_r}-\mathbb{F}_{c_r}$ with distortion $\tau$ if
  \[
    \left(1-\tau\right)\left\|\mathbf h\right\|_2^2
    \leq
    \left\|\mathbf A_\Omega \mathbf h\right\|_2^2
    \leq
    \left(1+\tau\right)\left\|\mathbf h\right\|_2^2,
    \qquad
    \mathbf h\in \mathbb{F}_{c_r}-\mathbb{F}_{c_r}.
  \]
\end{definition}

These conditions formalize the requirement that the measurement operator preserve the geometry of the relevant signal class. Empirical nondegeneracy ushers in the sample complexity guarantees via Christoffel sampling, which we then leverage for our uniform signal recovery guarantees. As empirical nondegeneracy implies the S-REC, we lay out the proposition below: 

\begin{proposition}[Empirical Nondegeneracy Implies S-REC]\label{prop: end implies srec}
  If $\mathbf A_\Omega$ is empirically nondegenerate on $\mathbb{F}_{c_r}-\mathbb{F}_{c_r}$ with distortion $\tau\in(0,1)$, then $
    \mathbf A_\Omega$ satisfies $ \mathrm{S\text{-}REC}\left(\mathbb{F}_{c_r},\sqrt{1-\tau},0\right)$.
\end{proposition}

The proof follows directly from Definition~\ref{def:empirical-nondegeneracy}, so we omit it. Unless explicitly stated otherwise, every high-probability statement below is with respect to the i.i.d.\ draw of $\Omega$ from $\mu_{c_s}$.

\subsection{Sampling-Recovery Mismatch Sample Complexity}
We now give explicit sample complexity conditions guaranteeing that the sampling matrix $\mathbf A_\Omega$ is stable on the recovery class under the assumption that the generator $G(\cdot, c_r)$ is a ReLU network, or more generally $L$-Lipschitz. For notational convenience, we define
$
  \underline{\mu}
  =
  \min_{i \in D}\mu_{c_s}\left(i\right)
$
to be the minimum probability assigned to a Fourier frequency under the sampling law $\mu_{c_s}$. We first state the sample complexity for ReLU networks, which we define below, based on~\cite[Definition 2]{berk2022fourier}.


\begin{definition}[Conditional Bias-Free Depth-$d$ ReLU Generator]
  \label{def:bias-free-relu}
  Let $\mathcal C$ be a conditioning space and let $d\geq 2$. We say that
  $G:\mathbb R^k\times\mathcal C\to\mathbb R^n$ is a conditional
  bias-free depth-$d$ feedforward ReLU generator with layer widths
   $k=k_0\leq k_1,\dots,k_{d-1}$ and $k_d=n$ if, for each fixed condition $c$, the induced map $G(\cdot, c):  \mathbb{R}^k \to \mathbb{R}^n$ has the form
  \[
    G(\mathbf z,c)
    =
    W_d(c)\sigma\!\left(
      W_{d-1}(c)\sigma\!\left(
        \cdots \sigma\!\left(W_1(c)\mathbf z\right)
      \right)
    \right),
  \]
  where $W_\ell(c)\in\mathbb R^{k_\ell\times k_{\ell-1}}$ for
  $\ell=1,\dots,d$, $\sigma(t)=\max\{t,0\}$ is applied coordinate-wise, and
  there are no additive bias vectors. We model conditioning abstractly by allowing the network weights to depend on $c$, so that each $c\in\mathcal C$ indexes a generator $G(\cdot,c)$.
\end{definition}

Appendix~\ref{app:relu-specialization} records the more general piecewise-linear theorem from which the following ReLU sample complexity bounds follows. 

\begin{theorem}[Sample Complexity for ReLU Networks]
  \label{thm:relu-srec}
  Let $c_r\in\mathcal C$ be the recovery prompt used to condition the generator $G(\cdot, c_r)$, and let $c_s\in\mathcal C$ be the sampling prompt used to draw $\Omega$ from $\mu_{c_s}$ and build $\mathbf A_\Omega$.
  Assume that $G\left(\cdot,c_r\right)$ is a bias-free depth-$d$ feedforward ReLU network in the sense of Definition~\ref{def:bias-free-relu} with layer widths
  $
    k = k_0 \leq k_1,\dots,k_{d-1},
    k_d = n
  $, with $d \geq 2$,
  and define
  \[
    \overline{k}_{c_r}
    =
    \left(
      \prod_{\ell=1}^{d-1} k_\ell
    \right)^{1/(d-1)}.
  \]
  Fix $0<\tau,\delta<1$. If
  \[
m \gtrsim \tau^{-2}\Lambda(c_r,c_r,c_s)\left[
  k(d-1)\left(1+\log\frac{\overline{k}_{c_r}}{k}\right)
  +
  k\left(1+\log\frac{1}{\tau\sqrt{\underline{\mu}}}\right)
  +
  \log\frac{1}{\delta}
\right],
\]
  then with probability at least $1-\delta$ over the draw of $\Omega$, $\mathbf A_\Omega$ satisfies
  $
  \mathrm{S\text{-}REC}
  \left(
    \mathbb{F}_{c_r},
    \sqrt{1-\tau},
    0
  \right)
  $.
\end{theorem}

We next record the sample complexity for stable recovery under the assumption that $G(\cdot, c_r)$ is $L$-Lipschitz, where proof is given in Appendix~\ref{app:lipschitz-specialization}.

\begin{theorem}[Sample Complexity for $L$-Lipschitz Networks]
  \label{thm:lipschitz-srec}
  Let $c_r\in\mathcal C$ be the recovery prompt used to condition the generator $G(\cdot, c_r)$, and let $c_s\in\mathcal C$ be the sampling prompt used to draw $\Omega$ from $\mu_{c_s}$ and build $\mathbf A_\Omega$.
  Assume that $G\left(\cdot,c_r\right)$ is $L$-Lipschitz on $B_2^k\left(R\right)$. Fix $0<\tau,\delta<1$ and $\xi>0$. If
  \[
    m
    \gtrsim
    \tau^{-2}
    \Lambda\left(c_r,c_r,c_s\right)
    \left[
      k\log\left(
        1+
        \frac{LR}
        {\xi\tau\sqrt{\underline{\mu}}}
      \right)
      +
      \log\left(\frac{1}{\delta}\right)
    \right],
  \]
  then with probability at least $1-\delta$ over the draw of $\Omega$, $\mathbf A_\Omega$ satisfies
  $
  \mathrm{S\text{-}REC}
  \left(
    \mathbb{F}_{c_r},
    \sqrt{1-\tau},
    \sqrt{1-\tau}\,\xi
  \right)
  $.
\end{theorem}

In both cases, conditioning affects sample complexity through alignment and coverage. Under sampling-recovery mismatch, the prompt compatibility factor $\Lambda(c_r,c_r,c_s)$ measures whether the sampling prompt $c_s$ places sufficient mass on the Fourier coordinates that distinguish signals in the recovery class $\mathbb F_{c_r}$, while $\underline{\mu}$ quantifies how close the sampling law is to neglecting coordinates altogether. If instead $c_s=c_r$, this penalty reduces to $\Lambda(c_r,c_r,c_r)=\kappa(\mathbb F_{c_r}-\mathbb F_{c_r})$, we recover the existing Christoffel- and coherence-based sample complexity results for generative compressed sensing~\cite{adcock2023cs4ml, berk2023modeladapted, plan2025denoising}. Under prompt mismatch, however, $\Lambda$ can grow and $\underline{\mu}$ can shrink, reflecting the increased number of samples required to capture poorly represented directions for stable recovery.


\subsection{Recovery Under the S-REC}\label{subsec:recovery}
Under the event that $\mathbf A_\Omega$ satisfies
$\mathrm{S\text{-}REC}\left(\mathbb{F}_{c_r},\gamma,q\right)$
for some $\gamma>0$ and $q\geq 0$, we can bound the recovery error. To do so, let $\widehat{\mathbf z}\in B_2^k(R)$ be an $\omega$-minimizer, meaning that
\begin{equation}\label{eq: omega minimizer}
  \left\|
  \mathbf y
  -
  \mathbf A_\Omega G\left(\widehat{\mathbf z},c_r\right)
  \right\|_2
  \le
  \inf_{\mathbf z\in B_2^k(R)}
  \left\|
  \mathbf y
  -
  \mathbf A_\Omega G\left(\mathbf z,c_r\right)
  \right\|_2
  +
  \omega,
\end{equation}
and likewise define $\widehat{\mathbf f}=G\left(\widehat{\mathbf z},c_r\right)$. Suppose the target signal $\mathbf f^*$ belongs to a different prompt-conditioned class $\mathbb F_{c_*}$. In this setting, $\Lambda(c_*,c_r,c_s)$ quantifies how the signal, recovery, and sampling prompt model classes interact, giving us the bound below:


\begin{theorem}[Prompt-Mismatched Recovery]
  \label{thm:prompt-mismatched-recovery}
 Let $c_r\in\mathcal C$ be the recovery prompt and let $\widehat{\mathbf z}$ denote an $\omega$-minimizer in the sense of Equation~\eqref{eq: omega minimizer} with estimator $\widehat{\mathbf f}=G(\widehat{\mathbf z},c_r)$. Let $c_s\in\mathcal C$ be the sampling prompt used to draw $\Omega$ from $\mu_{c_s}$ and build $\mathbf A_\Omega$, and let $c_* \in\mathcal C$ be the true-signal prompt so that $\mathbf f^* \in\mathbb{F}_{c_*}$.
  Suppose that $\mathbf A_\Omega$ satisfies $\mathrm{S\text{-}REC}\left(\mathbb{F}_{c_r},\gamma,q\right)$
  for some $\gamma>0$ and $q\geq 0$. Then
  \[
    \left\|\widehat{\mathbf f}-\mathbf f^*\right\|_2
    \le
    \left(
      1+
      \frac{2\sqrt{\Lambda\left(c_*,c_r,c_s\right)}}{\gamma}
    \right)
    \inf_{\mathbf f\in\mathbb{F}_{c_r}}
    \left\|\mathbf f^* - \mathbf f\right\|_2
    +
    \frac{2\left\|\mathbf e\right\|_2+\omega+q}{\gamma}.
  \]
\end{theorem}

Substituting Theorem~\ref{thm:relu-srec} yields $\gamma=\sqrt{1-\tau}$ and $q=0$, while substituting Theorem~\ref{thm:lipschitz-srec} yields $\gamma=\sqrt{1-\tau}$ and $q=\sqrt{1-\tau}\xi$, producing an additional additive error term of $\xi$. The recovery bound decomposes into an approximation error over the recovery class $\mathbb F_{c_r}$, an amplification factor governed by $\Lambda(c_*,c_r,c_s)$, and an additive term due to measurement noise, optimization error, and S-REC slack. If $\mathbf f^* \in \mathbb F_{c_r}$, this approximation term vanishes. Otherwise, the mismatch is controlled by the cross-class Christoffel function $K(\mathbb F_{c_*}-\mathbb F_{c_r})$ through $\Lambda(c_*,c_r,c_s)$, which measures how well $\mu_{c_s}$ covers Fourier directions separating the true-signal and recovery classes.

\section{Experiments}\label{sec: experiments}
In this section we examine both the prompt compatibility factor $\Lambda$ and its practical effect on conditional generative compressed sensing with Stable Diffusion 1.5 (SD15)~\cite{rombach2022high} through the Diffusers HuggingFace library~\cite{von-platen-etal-2022-diffusers}. Throughout, SD15 is used at resolution $512\times 512$ with $20$ DDIM~\cite{song2021denoising} denoising steps and classifier-free guidance (CFG)~\cite{ho2021classifierfree} of $7.5$ to produce signals $\mathbf{f} \in \mathbb{R}^{3 \times 512 \times 512}$ from a latent variable $\mathbf z \in \mathbb R^{4\times 64\times 64}$. In the reconstruction experiments we use the unweighted, unnormalized,  subsampled  Fourier operator
\[
  \mathbf A_{\Omega}
  =
  \tfrac{1}{\sqrt{m}} \mathbf P_{\Omega} \mathbf F,
  \qquad
  \mathbf y
  =
  \mathbf A_{\Omega} \mathbf f^*,
\]
so the prompt dependence enters through the selected Fourier mask $\Omega$. We do not introduce noise to the measurements. Since the images are multi-channel, we use a block Fourier sampling model: each sampled frequency in $\Omega$ is measured in all three color channels. We write $n=512\times 512$ for the size of the flattened image per channel, so the full ambient signal dimension is $3n$. To facilitate reproducibility, we provide the code for computing the empirical Christoffel sampling measures and running the reconstruction experiments.\footnote{Code available at: \href{https://github.com/alexdelise/ActiveConditionalGCS.git}{\texttt{github.com/alexdelise/ActiveConditionalGCS}}}


\subsection{Empirical Investigation into the Prompt Compatibility Factor  \texorpdfstring{$\Lambda$}{Lambda}}
We begin by studying how the prompt interaction factor changes under sampling-recovery mismatch $c_s \neq c_r$. To isolate this effect, we consider the prompt family
\begin{equation}\label{eq: prompt family}
    c_{\mathrm{uc}} = \texttt{""}, \qquad
    c_{\mathrm{db}} = \texttt{"daytime beach"}, \qquad
    c_{\mathrm{sb}} = \texttt{"sunset beach"}, \qquad
    c_{\mathrm{ca}} = \texttt{"cat"}.
\end{equation}
For each sampling prompt $c_s$, we use Algorithm~\ref{alg: ktilde} with $500$ iterations to estimate the empirical Christoffel function $\widetilde K$, construct the empirical sampling law $\widetilde\mu_{c_s}$, and compute $\widetilde{\Lambda}(c_r,c_r,c_s)$ following Equation~\eqref{eq: empirical lambda}. This lets us compare the prompt-matched regime $c_s=c_r$, an unconditioned baseline $c_s=c_{\mathrm{uc}}$, a mild semantic mismatch such as $c_{\mathrm{db}}$ versus $c_{\mathrm{sb}}$, and a severe mismatch induced by the unrelated class $c_{\mathrm{ca}}$. 

Figure~\ref{fig: distribution and heatmap} displays the resulting prompt-conditioned sampling distributions and the corresponding relative values of $\widetilde{\Lambda}$ to quantify how rapidly sample complexity deteriorates as the sensing prompt moves away from the recovery class. Immediately, one can see in Figure~\ref{fig: sampling_distributions} that the prompt-conditioned sampling distributions concentrate probability mass differently, but share a common bias toward low frequencies. The heatmap in Figure~\ref{fig: lambda_heatmap} makes the interaction between these differences readily apparent as the diagonal entries, corresponding to a matched sampling-recovery regime ($c_s = c_r$) minimize $\widetilde \Lambda$, with larger penalties occurring for semantically distant pairs in the sample-recovery mismatch regime, e.g. $c_{\mathrm{ca}}$ vs. $c_{\mathrm{uc}}$.

\begin{figure}
    \centering

    \begin{subfigure}{0.45\linewidth}
        \centering
        \includegraphics[width=\linewidth]{figures/sd15_lambda_figures/sd15_sampling_distribution_row_uc_sb_db_ca.pdf}
        \caption{}
        \label{fig: sampling_distributions}
    \end{subfigure}
    \hfill
    \begin{subfigure}{0.54\linewidth}
        \centering
        \includegraphics[width=\linewidth]{figures/sd15_lambda_figures/sd15_lambda_heatmap.pdf}
        \caption{}
        \label{fig: lambda_heatmap}
    \end{subfigure}

    \caption{
    \textbf{(a)} Empirical sampling distributions $\widetilde{\mu}_c$ induced by the Christoffel function for prompts $c \in \{c_{\mathrm{uc}}, c_{\mathrm{sb}}, c_{\mathrm{db}}, c_{\mathrm{ca}}\}$. Color indicates sampling probability over Fourier frequencies.
    %
    \textbf{(b)} Empirical prompt compatibility factor $\widetilde{\Lambda}(c_r,c_r,c_s)$ (rows: sampling, columns: recovery). Diagonal entries that represent $\widetilde \kappa$ are minimal, while off-diagonal entries indicate the sampling-recovery mismatch penalty.
    }
    \label{fig: distribution and heatmap}
\end{figure}


%%%%%%%%%%%%%%%%%%%%%%%%%%%%%%%%%%%%%%%%%%%%%%%%%%%%%%%%%%%%%%%%%%%%%%%%%%%%%%
%%%%%%%%%%%%%%%%%%%%%%%%%%%%%%%%%%%%%%%%%%%%%%%%%%%%%%%%%%%%%%%%%%%%%%%%%%%%%%
%%%%%%%%%%%%%%%%%%%%%%%%%%%%%%%%%%%%%%%%%%%%%%%%%%%%%%%%%%%%%%%%%%%%%%%%%%%%%%
\subsection{Conditional Generative Compressed Sensing}
We next connect these prompt-conditioned sampling laws to actual reconstruction performance. For each pair of sampling prompt $c_s$ and recovery prompt $c_r$ in Equation~\eqref{eq: prompt family}, we reconstruct the same target image across sampling rates $m / n$ using $5$ independently drawn masks per setting. For each mask, we always sample the $0$ frequency, a standard assumption in compressive imaging~\cite{adcock2021compressive}, and then sample the remaining $m-1$ frequencies via $\widetilde \mu_{c_s}$ without replacement.

Given an image $\mathbf{f}^*$ and measurements $\mathbf y = \mathbf A_{\Omega} \mathbf f^*$, we reconstruct under prompt $c_r$ by seeking $\widehat{\mathbf f}=G(\widehat{\mathbf z},c_r)$, where $\widehat{\mathbf z}$ is obtained from the least-squares problem
\begin{equation}\label{eq: least squares}
  \widehat{\mathbf z}
  \in
  \arg\min_{\mathbf z}
  \frac{1}{2 \cdot 3m}
  \left\|
    \mathbf A_{\Omega} G(\mathbf z,c_r) - \mathbf y
  \right\|_2^2.
\end{equation}


We initialize the generator from the zero-filled inverse Fourier transform of our subsampled measurement
$
  \mathbf f_{\mathrm{zf}}
  =
  \tfrac{m}{n}\mathbf A_{\Omega}^H \mathbf y,
$
where $H$ denotes the Hermitian adjoint, by encoding $\mathbf f_{\mathrm{zf}}$ into latent space and then progressively adding noise to obtain $\mathbf z_{\mathrm{init}}$. Intuitively, this initialization anchors the reconstruction to the measured image content, then lifts it to the noise scale used at the first reverse diffusion step. We treat the composition of the $20$-step DDIM denoising process and final VAE decode as the generative model $G(\cdot,c_r)$. After sampling $\Omega \sim \widetilde{\mu}_{c_s}$, we solve Equation~\eqref{eq: least squares} by optimizing using Adam~\cite{kingma2015adam} in PyTorch with base learning rate $10^{-2}$, cosine decay, gradient clipping at $1.0$, and early stopping after $35$ stagnant optimization steps, for at most $500$ steps. All experiments were performed across five NVIDIA RTX A5000 GPUs with 24GB VRAM for a total runtime of approximately ten days. To see the ground truth images as well as example reconstructions under each sampling-recovery prompt pair in the experiments, we refer the reader to Appendix~\ref{app: recovered img comp}. Solid lines in figures denote mean reconstruction metrics across the five independent trials at each sampling percentage, with shading denoting $95\%$ confidence intervals computed from the standard error of the mean.




%%%%%%%%%%%%%%%%%%%%%%%%%%%%%%%%%%%%%%%%%%%%%%%%%%%%%%%%%%%%%%%%%%%%%%%%%%%%%%
%%%%%%%%%%%%%%%%%%%%%%%%%%%%%%%%%%%%%%%%%%%%%%%%%%%%%%%%%%%%%%%%%%%%%%%%%%%%%%
%%%%%%%%%%%%%%%%%%%%%%%%%%%%%%%%%%%%%%%%%%%%%%%%%%%%%%%%%%%%%%%%%%%%%%%%%%%%%%
\subsubsection{In-Range Prompt-Mismatched Image Recovery}\label{subsec: in-range diff prompt}
In this scenario, the target image we seek to reconstruct is generated by SD15, using the prompt $c_* = \texttt{"sunset over a sandy coast"}$, so
$
  \mathbf f^*
  =
  G(\mathbf z^*, c_*)
$
for some latent vector $\mathbf z^*$. Since $c_*$ does not appear in Equation~\eqref{eq: prompt family}, we consider this an \emph{in-range prompt-mismatched image recovery} task. Figure~\ref{fig: in-range prompt-mismatch recon metrics} demonstrates how sampling from $\widetilde \mu_{c_\mathrm{sb}}$ improves reconstruction performance compared to sampling from the other distributions, as is expected due to the semantic similarity of $c_\mathrm{sb}$ to $c_*$. Under unconditioned recovery ($c_r = c_{\mathrm{uc}}$), sampling from $\widetilde \mu_{c_\mathrm{sb}}$ improves PSNR relative to the average of the other sampling laws at every tested sampling percentage, with gains ranging from $0.62$ to $2.96$ dB, and with corresponding SSIM gains ranging from $0.0034$ to $0.0469$. Since we see best performance across all sampling distributions with unconditioned reconstruction, this suggests that an unconditioned generator may be more flexible during reconstruction, especially with CFG $7.5$. All sampling-reconstruction prompt pairs far outperform the zero-filled image baseline in terms of mean SSIM, indicating that the generator provides a useful image prior even under prompt mismatch.


\begin{figure}
    \centering
    \includegraphics[width=1.0\linewidth]{figures/sunset_recovery_final/diffusion_backprop_cs_combined_metrics_by_recovery_prompt.pdf}
    \caption{Reconstruction quality across sampling percentages by sampling distribution for the in-range prompt-mismatched experiment.}
    \label{fig: in-range prompt-mismatch recon metrics}
\end{figure}




%%%%%%%%%%%%%%%%%%%%%%%%%%%%%%%%%%%%%%%%%%%%%%%%%%%%%%%%%%%%%%%%%%%%%%%%%%%%%%
%%%%%%%%%%%%%%%%%%%%%%%%%%%%%%%%%%%%%%%%%%%%%%%%%%%%%%%%%%%%%%%%%%%%%%%%%%%%%%
%%%%%%%%%%%%%%%%%%%%%%%%%%%%%%%%%%%%%%%%%%%%%%%%%%%%%%%%%%%%%%%%%%%%%%%%%%%%%%
\subsubsection{Out-of-Range Image Recovery}\label{subsec: out of range}
For the \emph{out-of-range} image recovery task, we follow the same procedure, although this time the sensed image is obtained from~\cite{magnificSunsetBeachPalmTree}. From Figure~\ref{fig: out-of-range recon metrics}, we observe that the out-of-range reconstructions are worse than the in-range prompt-mismatched reconstructions in PSNR, while SSIM is comparable and in some cases slightly higher. Paired against the in-range prompt-mismatched experiment the out-of-range reconstructions have $3.50$ dB lower mean PSNR, while mean SSIM is slightly higher by $0.0094$. As before, we observe that reconstruction over $c_r = c_{\mathrm{uc}}$ provides the best results, however, the gain achieved from sampling from $\widetilde\mu_{c_\mathrm{sb}}$ is diminished here. Under unconditioned recovery, the PSNR gain from $\widetilde\mu_{c_\mathrm{sb}}$ relative to the average of the other sampling laws drops from $1.34$ dB in the in-range prompt-mismatched experiment to $0.13$ dB in the out-of-range experiment. We attribute this to the fact that $\widetilde\mu_{c_\mathrm{sb}}$ is constructed from images within the range of the generator, and thus may not emphasize Fourier frequencies most relevant for reconstructing signals that lie outside this range.


\begin{figure}
    \centering
    \includegraphics[width=1.0\linewidth]{figures/sunset_ood_recovery_final/diffusion_backprop_cs_combined_metrics_by_recovery_prompt.pdf}
    \caption{Reconstruction quality across sampling percentages by sampling distribution for the out-of-range experiment.}
    \label{fig: out-of-range recon metrics}
\end{figure}



\subsection{Results Discussion}
Taken together, the experiments support the value of prompt-conditioned sampling while also highlighting limitations of the current experimental pipeline.
The empirical prompt compatibility factor $\widetilde{\Lambda}$ values indicate that sampling-recovery mismatch is indeed reflected in the prompt-conditioned sampling laws, with matched prompt pairs producing smaller compatibility penalties. This aligns with the theory, where $\Lambda(c_r,c_r,c_s)$ controls the sample complexity  for stable recovery. However, the reconstruction results do not perfectly follow the ordering suggested by this factor. In particular, reconstruction with $c_r=c_{\mathrm{uc}}$ consistently performs best when sampling from all distributions, even over reconstructing with $c_r = c_{\mathrm{sb}}$. 
We attribute this to three main effects. First, the experiments use an unweighted least-squares objective rather than the weighted objective analyzed in Section~\ref{sec: theory}, so rare but high-leverage frequencies are not amplified according to $1/\mu_{c_s}(i)$ and may be underfit in practice. Second, the sampling laws are based on finite Monte Carlo estimates of the Christoffel function, which can underestimate worst-case Fourier directions. Since $\Lambda$ is a max-ratio quantity, small errors in the tails of $\widetilde{\mu}_{c_s}$ can produce large practical mismatch penalties. Third, the latent optimization problem is highly nonconvex and optimizer-dependent. The unconditional model may provide a less restrictive and more easily optimized reconstruction class, especially under high CFG, even if it is not theoretically the best-matched prompt class. For further exploration on the role of CFG in sampling and reconstruction, see Appendix~\ref{app: cfg ablation}.

\section{Conclusion}\label{sec: conclusion}

We develop an active learning framework for conditional generative compressed sensing in which prompt-conditioned generative models are used for signal recovery and constructing sampling distributions for subsampled Fourier measurements. We introduce a prompt compatibility factor $\Lambda$ that quantifies alignment between the sampling, recovery, and ground-truth prompt-conditioned signal classes, and demonstrate how it governs sample complexity under ReLU and Lipschitz architecture assumptions and bounds reconstruction error. Experiments with Stable Diffusion show that prompt conditioning can substantially alter the induced Christoffel sampling distributions and can influence reconstruction performance under subsampled Fourier measurements.

Our work suggests several directions for future research. Since the prompt compatibility factor $\Lambda$ is defined through a max-divergence quantity, future work could investigate KL- or R\'{e}nyi divergence-versions of $\Lambda$ and the resulting sample complexity for (nonuniform) recovery. Also, our theory applies to idealized generator classes and is not yet well suited to modern conditional generative models. Although a model such as SD15 can be viewed as inducing a Lipschitz map on compact latent sets after fixing the prompt embedding, sampler, denoising schedule, and CFG scale, this remains an idealized abstraction of practical recovery. The effective Lipschitz constant may be large and sensitive to modeling choices such as CFG scale, suggesting that such parameters, together with prompts, should be studied as design variables for the sampling and recovery classes.






%\section*{References}
%\clearpage
\bibliographystyle{unsrtnat}
\bibliography{references}


%%%%%%%%%%%%%%%%%%%%%%%%%%%%%%%%%%%%%%%%%%%%%%%%%%%%%%%%%%%%

\appendix
%\newpage
\newpage\section{Christoffel Sampling for Machine Learning (CS4ML)}\label{app:cs4ml}

In this appendix, we fully outline the abstract CS4ML framework in~\cite{adcock2023cs4ml} we harness for Christoffel sampling and bridge their objects to our weighted subsampled conditional Fourier model.

\subsection{Objects, Definitions, and Assumptions}
We begin with an underlying probability space $(\Omega,\mathcal{F},\mathbb{P})$, a separable Hilbert space $\mathbb{X}$, and a normed vector subspace $\mathbb{X}_0 \subseteq \mathbb{X}$, termed the \emph{object space}. Our goal is to learn an object $\mathbf{f}^* \in \mathbb{X}_0$ from training data. We consider the measurement domain $(D,\mathcal{A},\rho)$, the semi-inner-product measurement space $\mathbb{Y}$, and the sampling operator $L : D \to \mathcal B(\mathbb{X}_0,\mathbb{Y})$, where $\mathcal B(\mathbb{X}_0,\mathbb{Y})$ denotes the space of bounded linear operators from $\mathbb{X}_0$ to $\mathbb{Y}$. Thus, for each $i \in D$, the quantity $L(i)(\mathbf{f}^*) \in \mathbb{Y}$ represents the measurement of $\mathbf{f}^*$ at parameter value $i$ (where we employ the classical active learning assumption that we may query $L$ for any $i \in D$). We also fix an \emph{approximation space} $\mathbb{F} \subseteq \mathbb{X}_0$ within which reconstruction is sought, together with a sampling measure $\mu$ on $(D,\mathcal{A})$ that is absolutely continuous with respect to $\rho$, with Radon--Nikodym derivative $\nu = d\mu/d\rho > 0$.


To ensure that the sampling operators are sufficiently expressive to preserve the $\mathbb{X}$-norm we make the following assumption.

\begin{assumption}[Nondegeneracy of the sampling operators]
    There exist constants $0<\alpha\leq \beta<\infty$ such that
    \[
      \alpha\left\|\mathbf{f}\right\|_{\mathbb{X}}^2
      \leq
      \int_{D}
      \left\|L(i)(f)\right\|_{\mathbb{Y}}^2
      \, d\rho(i)
      \leq
      \beta\left\|\mathbf{f}\right\|_{\mathbb{X}}^2,
      \qquad
      \mathbf{f}\in\mathbb{X}_0,
      \qquad 
      i \in D.
    \]
\end{assumption}
Later, we will define the discrete analogue of nondegeneracy in which we hope to approximately preserve norms.

For a fixed space $\mathbb{F}$, the generalized Christoffel function and its related quantity are defined as
\[
  K(\mathbb{F})(i)
  =
  \sup_{\mathbf{f}\in \mathbb{F}\setminus\{0\}}
  \frac{\left\|L(i)(\mathbf{f})\right\|_{\mathbb{Y}}^2}
       {\left\|\mathbf{f}\right\|_{\mathbb{X}}^2},
  \qquad
  \kappa(\mathbb{F})
  =
  \int_{D} K(\mathbb{F})(i)\, d\rho(i),
  \qquad 
  i \in D.
\]
The corresponding optimal sampling density is
\[
  \nu^*(i)
  =
  \frac{K(\mathbb{F})(i)}{\kappa(\mathbb{F})},
  \qquad 
  i \in D,
\]
 given by~\cite[Lemma~4.6]{adcock2023cs4ml}. Equivalently, the optimal sampling measure $\mu^*$ is the measure with density $\nu^*$ with respect to $\rho$. Finally, the sampling seminorm for a density $\nu$ is given by
\[
  \snorm{\mathbf{f}}{\nu}
  =
  \left(
    \operatorname*{ess\,sup}_{i \sim \rho}
    \frac{\left\|L(i)(\mathbf{f})\right\|_{\mathbb{Y}}^2}{\nu(i)}
  \right)^{1/2}.
\]

\subsection{Specialization to Weighted Subsampled Fourier Measurements}

For our paper, the aforementioned objects above collapse to the following for subsampled Fourier measurements:
\[
  \mathbb{X}=\mathbb{X}_0=\mathbb{R}^n,
  \qquad
  \left\|\cdot\right\|_{\mathbb{X}}=\left\|\cdot\right\|_2,
  \qquad
  D=\{1,\dots,n\},
\]
\[
  \rho=\text{counting measure on }D,
  \qquad
  \mathbb{Y}=\mathbb{C}
  \text{ with }
  \left\|z\right\|_{\mathbb{Y}}=|z|,
  \qquad
  L(i)(\mathbf f)=\mathbf P_i\mathbf F\mathbf f.
\]
Since $\mathbf{F}$ is the unitary discrete Fourier transform, Parseval's theorem gives
\[
  \sum_{i=1}^n |L(i)(\mathbf f)|^2
  =
  \sum_{i=1}^n \left|\mathbf P_i\mathbf F\mathbf f\right|^2
  =
  \left\|\mathbf f\right\|_2^2,
  \qquad
  \mathbf f\in\mathbb{R}^n,
\]
so the nondegeneracy constants are $\alpha=\beta=1$. Also, as $D$ is finite, we make the following assumption:

\begin{assumption}
  \label{assump: full-support-general}
  For every sampling prompt $c\in\mathcal C$, the associated sampling law
  $\mu_c$ is a probability vector on $D=\{1,\dots,n\}$ satisfying
  \[
    \mu_c(i)>0,
    \qquad i\in D.
  \]
  Equivalently, if $\nu_c=d\mu_c/d\rho$ denotes the density of $\mu_c$
  with respect to counting measure $\rho$, then
  \[
    \nu_c(i)>0,
    \qquad i\in D.
  \]
\end{assumption}

\subsection{Prompt-Conditioned Discrete Analogues}

We now return to the prompt-conditioned recovery setting from the main text by fixing a recovery prompt $c_r$ and applying the CS4ML machinery to the secant class $\mathbb{F}_{c_r}-\mathbb{F}_{c_r}$. Under this specialization, the optimal density $\nu^*$ becomes the Christoffel sampling law $\mu_{c_s}$, and thus
\[
  \operatorname*{ess\,sup}_{i\sim\rho}
  \frac{K(\mathbb{F}_{c_r}-\mathbb{F}_{c_r})(i)}{\nu^*(i)}
  =
  \max_{i \in D}
  \frac{K(\mathbb{F}_{c_r}-\mathbb{F}_{c_r})(i)}{\mu_{c_s}(i)}
  =
  \Lambda(c_r,c_r,c_s).
\]
The optimal sampling seminorm reduces to
\[
  \snorm{\mathbf f}{\nu^*}
  =
   \left(
    \operatorname*{ess\,sup}_{i \sim \rho}
    \frac{\left\|L(i)(\mathbf{f})\right\|_{\mathbb{Y}}^2}{\nu(i)}
  \right)^{1/2} 
  = 
  \left(
    \max_{i \in D}
    \frac{|\mathbf P_i\mathbf F\mathbf f|^2}{\mu_{c_s}(i)}
  \right)^{1/2}
  =
  \snorm{\mathbf f}{\mu_{c_s}},
\]
and we have that
\[
  \frac{1}{m}\sum_{j=1}^m
  \frac{|L(I_j)(\mathbf f)|^2}{\mu_{c_s}(I_j)}
  =
  \frac{1}{m}\sum_{j=1}^m
  \frac{|\mathbf P_{I_j}\mathbf F\mathbf f|^2}{\mu_{c_s}(I_j)}
  =
  \left\|\mathbf A_\Omega \mathbf f\right\|_2^2.
\]

We also write
\[
  \mathbb B_{c_r}
  =
  \left\{
    \mathbf h/\left\|\mathbf h\right\|_2
    :
    \mathbf h\in
    \left(\mathbb F_{c_r}-\mathbb F_{c_r}\right)\setminus\{\mathbf 0\}
  \right\}
\]
for the normalized recovery secant set.

 Under this specialization, empirical nondegeneracy is given by the following, which we repeat for completeness:

\begin{definition}[Empirical Nondegeneracy]
  Let $0<\tau<1$. We say that $\mathbf A_\Omega$ is \emph{empirically nondegenerate} on $\mathbb{F}_{c_r}-\mathbb{F}_{c_r}$ with distortion $\tau$ if
  \[
    \left(1-\tau\right)\left\|\mathbf f\right\|_2^2
    \leq
    \left\|\mathbf A_\Omega \mathbf f\right\|_2^2
    \leq
    \left(1+\tau\right)\left\|\mathbf f\right\|_2^2,
    \qquad
    \mathbf f\in \mathbb{F}_{c_r}-\mathbb{F}_{c_r}.
  \]
\end{definition}

Finally, the theorem below is the discrete subsampled Fourier measurement sample-complexity version of \cite[Theorem~4.2]{adcock2023cs4ml}, which will be particularly useful when proving Theorem~\ref{thm:relu-srec}.

\begin{theorem}[Empirical Nondegeneracy and S-REC Under Christoffel Sampling]
  \label{thm:christoffel-srec}
  Fix $0<\tau,\delta<1$, and let $\mathcal N\subseteq \mathbb{B}_{c_r}$ be an $\eta$-net for $\mathbb{B}_{c_r}$ in the sampling seminorm $\snorm{\cdot}{\mu_{c_s}}$ with $0<\eta\leq \tfrac{\tau}{8}$. If
  \[
    m \gtrsim \tau^{-2}\Lambda\left(c_r,c_r,c_s\right)\log\left(\frac{2\left|\mathcal N\right|}{\delta}\right),
  \]
  then with probability at least $1-\delta$ over the draw of $\Omega$, $\mathbf A_\Omega$ is empirically nondegenerate on $\mathbb{F}_{c_r}-\mathbb{F}_{c_r}$ with distortion $\tau$, and hence satisfies
  $
  \mathrm{S\text{-}REC}\left(\mathbb{F}_{c_r},\sqrt{1-\tau},0\right).
  $
\end{theorem}

We defer the proof to Appendix~\ref{app:theory-proofs}.

\newpage \section{Numerical Approximation to the Generalized Christoffel Function}\label{app: further experiment details}
In practice, the Christoffel function $K(\mathbb F_{c_s}-\mathbb F_{c_s})$ appearing in the sampling law in Equation~\eqref{eq: christoffel sampling law} is not computable exactly when the prompt-conditioned class $\mathbb F_{c_s}$ is specified as the range of a deep generator $G(\cdot,c_s)$. Indeed, evaluating $K(\mathbb F_{c_s}-\mathbb F_{c_s})(i)$ requires taking a supremum over all nonzero secants of $\mathbb F_{c_s}$, equivalently over pairs of latent vectors passed through $G(\cdot,c_s)$, which amounts to a global optimization problem that is intractable at scale. We therefore use Algorithm~\ref{alg: ktilde}, adapted from~\cite[Algorithm~1]{adcock2023cs4ml}, to construct an empirical estimator $\widetilde K$ of $K(\mathbb F_{c_s}-\mathbb F_{c_s})$. This estimator is straightforward to compute in practice and is intended to approximate, and in the large-sample limit recover, the true Christoffel function.

With $\widetilde K$ in hand, we define the empirical analogue of the Christoffel sampling law in Equation~\eqref{eq: christoffel sampling law} by
\begin{equation}\label{eq: empirical sampling law}
  \widetilde \mu_{c_s}(i)
  =
  \frac{\widetilde K(i)}{\widetilde \kappa_{c_s}},
  \qquad
  \widetilde \kappa_{c_s}
  =
  \sum_{j=1}^{n} \widetilde K(j),
  \qquad
  i \in D.
\end{equation}
This yields a computable approximation to the ideal sampling law $\mu_{c_s}$, and we then draw the measurement indices in $\Omega$ independently from $\widetilde \mu_{c_s}$. Finally, we can compute an empirical approximation to the prompt compatibility factor using the formula 
\begin{equation}\label{eq: empirical lambda}
 \widetilde\Lambda(c_1,c_2,c_3)
    =
    \max_{i \in D}
    \frac{\widetilde K(\mathbb{F}_{c_1}-\mathbb{F}_{c_2})(i)}{\widetilde \mu_{c_3}(i)}.
\end{equation}

\begin{algorithm}
\caption{Compute $\widetilde{K}$}
\label{alg: ktilde}
\begin{algorithmic}[1]
\REQUIRE Generator $G$, sampling prompt $c_s$, latent distribution $\pi_z$, number of trials $t$
\STATE Initialize $\widetilde K(i)\gets 0$ for $i=1,\dots,n$
\FOR{$\ell=1$ to $t$}
    \STATE Sample latent vectors $\mathbf z_1,\mathbf z_2 \sim \pi_z$
    \STATE Generate $\mathbf f_1 \gets G(\mathbf z_1,c_s)$ and $\mathbf f_2 \gets G(\mathbf z_2,c_s)$
    \STATE Set $\mathbf h \gets \mathbf f_1-\mathbf f_2$ and $\epsilon \gets \left\|\mathbf h\right\|_2$
    \IF{$\epsilon = 0$}
        \STATE \textbf{continue}
    \ENDIF
    \STATE Compute the discrete Fourier transform $\mathbf y \gets \mathbf F \mathbf h$
    \FOR{$i=1$ to $n$}
        \STATE $a_i \gets |y_i|^2 / \epsilon^2$
        \STATE $\widetilde K(i) \gets \max\{\widetilde K(i), a_i\}$
    \ENDFOR
\ENDFOR
\STATE Set $\widetilde \mu_{c_s}(i) \gets \widetilde K(i) / \sum_{j=1}^n \widetilde K(j)$ for $i=1,\dots,n$
\end{algorithmic}
\end{algorithm}


\newpage \section{In-Range Prompt-Matched Image Recovery Experiment}\label{app: in range prompt match exp}
We perform another in-range image recovery experiment where the target image is itself generated by SD15, however this time we generate with the prompt $c_* = c_r =  \texttt{"sunset beach"}$, so
$
  \mathbf f^*
  =
  G(\mathbf z^*, c_*)
$
for some latent vector $\mathbf z^*$. Since $c_*$ appears in the prompt family in Equation~\eqref{eq: prompt family}, we consider this an \emph{in-range prompt-matched} image recovery task.  After sampling $\Omega \sim \widetilde{\mu}_{c_s}$, we solve Equation~\eqref{eq: least squares} by optimizing using Adam with a learning rate of $10^{-2}$, gradient clipping at $1.0$,  and early stopping after 25 stagnant optimization steps, for at most $500$ steps. As before, we recover over all combinations of $c_s$ and $c_r$, where each prompt is given from the prompt family in Equation~\eqref{eq: prompt family} for five trials at each sampling ratio. Appendix~\ref{app: recovered img comp} records reconstruction examples, as well as the ground truth image for this experiment.

Figure~\ref{fig: in-range prompt-match recon metrics} depicts image reconstruction metrics under severe undersampling across sampling distributions and image reconstruction prompts. In the scenario where we both sample and recover over $c_s = c_r = c_*$, one would expect reconstruction performance to be optimal, however in practice we observe that this is not the case. For example, at sampling ratio $0.00125$, the matched pair $c_s=c_r=c_\mathrm{sb}$ ranks sixth out of the sixteen sampling-recovery prompt pairs in both PSNR and SSIM, while sampling from $\widetilde\mu_{c_\mathrm{sb}}$ and reconstructing with $c_r=c_\mathrm{uc}$ gives the best mean PSNR and SSIM. As mentioned before, this is likely due to the high CFG value of 7.5 set for the generator, which may induce too much bias during reconstruction. Nevertheless, we see improved reconstruction when sampling from $\widetilde \mu_{c_\mathrm{sb}}$ and reconstructing with prompt $c_{\mathrm{uc}}$, in which we also see a tight confidence interval, even at extreme undersampling. This suggests that $\widetilde \mu_{c_\mathrm{sb}}$ places greater emphasis on Fourier frequencies that are important for accurately reconstructing images generated with prompt $c_r = c_\mathrm{sb}$ as expected by its construction.



\begin{figure}[htp]
    \centering
    \includegraphics[width=1.0\linewidth]{figures/sunset_recovery_full/diffusion_backprop_cs_combined_metrics_by_recovery_prompt.pdf}
    \caption{Reconstruction quality across sampling percentages by sampling distribution for the in-range prompt-matched experiment.}
    \label{fig: in-range prompt-match recon metrics}
\end{figure}
\newpage\section{CFG Ablation}\label{app: cfg ablation}
For diffusion generative image models, CFG controls how strongly the generated samples adhere to conditioning, in our case a text prompt. Higher CFG values produce images with better semantic alignment to the text prompt, whereas lower CFG values allow the model more diversity and natural variability~\cite{ho2021classifierfree}. For SD15, a CFG scale of $1$ corresponds to the baseline conditional generation process without additional CFG.

In the context of our framework, changing CFG effectively changes the prompt-conditioned generator itself: a low CFG value allows the reconstruction class to remain relatively broad, while a high CFG value restricts the generator more strongly toward images that match the prompt. Thus, CFG can affect both the sampling side and the recovery side of the problem. On the sampling side, it can change the geometry of the prompt-conditioned image class and therefore alter the estimated Christoffel sampling distribution. On the recovery side, it can change how much flexibility the optimizer has to match the observed Fourier measurements. We therefore include a small CFG ablation to study the tradeoff between stronger prompt adherence and reconstruction fidelity. 

\subsection{Sampling Distributions}
We begin by exploring how changing the CFG scale affects the prompt-conditioned Christoffel sampling distributions. Figure~\ref{fig: ktilde ablation} displays the sampling distributions obtained by estimating the empirical Christoffel function using Algorithm~\ref{alg: ktilde} with $500$ iterations under varying CFG scales, and then applying Equation~\ref{eq: empirical sampling law} to obtain $\widetilde \mu_c$. We observe that as the CFG scale increases, the sampling distributions become less diffuse and place more mass on a smaller set of Fourier frequencies. This is consistent with the role of CFG: stronger guidance restricts the generator more tightly to images aligned with the prompt, reducing variability within the prompt-conditioned class. Consequently, the empirical secants used to estimate the Christoffel function vary along fewer dominant directions, making the most informative Fourier frequencies easier to identify. Conversely, at lower CFG scales, the generator produces a broader and more variable collection of images, so the resulting secant energy is spread across a wider range of frequencies and the induced sampling law is more diffuse.


\begin{figure}[htp]
    \centering
    \includegraphics[width=1.0\linewidth]{figures/sd15_lambda_figures/sd15_sampling_distribution_cfg1_cfg3_cfg5_cfg7p5_grid.pdf}
    \caption{ Empirical sampling distributions $\widetilde{\mu}_c$ induced by the Christoffel function for prompts $c \in \{ c_{\mathrm{sb}}, c_{\mathrm{db}}, c_{\mathrm{ca}}\}$ across varying CFG scales. Color indicates sampling probability over Fourier frequencies.}
    \label{fig: ktilde ablation}
\end{figure}

\subsection{Image Recovery}
We next perform an ablation on CFG scale to study its effect on image recovery under the reconstruction prompt $c_r = c_{\mathrm{sb}} = \texttt{"sunset beach"}$, including unconditioned reconstruction ($c_r = c_\mathrm{uc}$) as a baseline. For each trial, we solve Equation~\eqref{eq: least squares} using Adam with learning rate $10^{-2}$, gradient clipping at $1.0$, and early stopping after $25$ stagnant optimization steps, for at most $500$ steps. For each sampling-recovery prompt pair and sampling percentage, we perform two reconstruction trials, except at CFG $7.5$, where we reuse the five trials from the primary experiments. All sampling was performed with respect to the original prompt-conditioned sampling distributions made with CFG $7.5$ in Figure~\ref{fig: sampling_distributions}.

Figures~\ref{fig: ablation in-range prompt-mismatched}--\ref{fig: ablation in-range prompt-matched} demonstrate how in all recovery scenarios, increasing CFG tends to decrease reconstruction performance, with lower CFG values improving reconstruction relative to the high-CFG setting used in the primary experiments. Averaged over sampling distributions and sampling percentages, recovering over $c_r = c_\mathrm{sb}$ with CFG $1$ improves over CFG $7.5$ by $2.31$ dB in the in-range prompt-mismatched experiment, $3.89$ dB in the in-range prompt-matched experiment, and $2.99$ dB in the out-of-range experiment; the corresponding SSIM gains are $0.0208$, $0.0346$, and $0.0328$. Across all three experiments, CFG $1$ outperforms CFG $7.5$ in all $60$ sampling-distribution/sampling-percentage cells for both PSNR and SSIM. In particular, we observe that recovering over $c_r = c_\mathrm{sb}$ with CFG $1$ matches or slightly outperforms recovering over $c_r = c_\mathrm{uc}$ on average, with mean gains of $0.123$ dB and $0.0012$ SSIM across all three experiments, although this advantage is not uniform across every sampling percentage and sampling law. This suggests that the weaker performance of prompt-conditioned recovery in the primary experiments in Sections~\ref{subsec: in-range diff prompt} and \ref{subsec: out of range} is not due to the prompt itself, but rather to the restrictive effect of large CFG values during latent optimization. As CFG increases, the generator is forced more strongly toward images semantically aligned with the prompt, which can reduce its ability to match the observed Fourier measurements. Lower CFG values therefore appear to provide a more flexible effective recovery class, allowing the optimizer to better balance prompt consistency with measurement fidelity.

Since all sampling masks in this ablation are drawn from distributions constructed at CFG $7.5$, these results primarily isolate the effect of CFG during recovery rather than during sampling. Taken together, this indicates that CFG is an important reconstruction hyperparameter in prompt-conditioned generative compressed sensing with diffusion models, and that future experiments should jointly tune the CFG scale used to construct $\widetilde \mu_c$ and the CFG scale used during recovery.


\begin{figure}
    \centering
    \includegraphics[width=1.0\linewidth]{figures/sunset_recovery_final_cfg_ablation/sd15_cfg_ablation_metric_curves.pdf}
    \caption{Reconstruction performance versus sampling percentage for in-range prompt-mismatched recovery under varying CFG scales, including unconditioned baseline.}
    \label{fig: ablation in-range prompt-mismatched}
\end{figure}


\begin{figure}
    \centering
    \includegraphics[width=1.0\linewidth]{figures/sunset_ood_recovery_final_cfg_ablation/sd15_cfg_ablation_metric_curves.pdf}
    \caption{Reconstruction performance versus sampling percentage for out-of-range recovery under varying CFG scales, including unconditioned baseline.}
    \label{fig: ablation out-of-range}
\end{figure}


\begin{figure}
    \centering
    \includegraphics[width=1.0\linewidth]{figures/sunset_recovery_full_cfg_ablation/sd15_cfg_ablation_metric_curves.pdf}
    \caption{Reconstruction performance versus sampling percentage for in-range prompt-matched recovery under varying CFG scales, including unconditioned baseline.}
    \label{fig: ablation in-range prompt-matched}
\end{figure}


\clearpage \section{Proofs}\label{app:theory-proofs}

In this appendix we collect the proofs for the arguments used throughout the paper. As a reminder, we redefine the following items. Let 
\[
  \mathbb B_{c_r}
  =
  \left\{
    \frac{\mathbf h}{\left\|\mathbf h\right\|_2}
    :
    \mathbf h\in
    \left(\mathbb F_{c_r}-\mathbb F_{c_r}\right)
    \setminus\{\mathbf 0\}
  \right\}
\]
denote the normalized secant set. For $\xi>0$, we also define the truncated normalized secant set by
\[
  \mathbb{B}_{c_r,\xi}
  =
  \left\{
    \frac{\mathbf h}{\left\|\mathbf h\right\|_2}
    :
    \mathbf h\in
    \left(\mathbb{F}_{c_r}-\mathbb{F}_{c_r}\right)\setminus\left\{\mathbf 0\right\},
    \ \left\|\mathbf h\right\|_2\geq \xi
  \right\},
\]
which will come into play when proving Lipschitz guarantees. Finally, we define
\[
\snorm{\mathbf f}{\mu_{c_s}} = 
\left(
    \max_{i \in D}
    \frac{|\mathbf P_i\mathbf F\mathbf f|^2}{\mu_{c_s}(i)}
  \right)^{1/2}
  \quad \text{and} \qquad 
  \underline{\mu}
  =
  \min_{i \in D}\mu_{c_s}\left(i\right),
\]
which are well-defined under Assumption~\ref{assump: full support}.

\subsection{Proofs of the CS4ML Secant-Stability Statements}\label{app:end-to-srec}

We begin with a quantitative secant-stability estimate in the sampling seminorm, from which the empirical-nondegeneracy and zero-slack S-REC statement under Christoffel sampling in Theorem~\ref{thm:christoffel-srec} follows.


\begin{theorem}
  \label{thm:secant-stability}
    Fix $0<\varepsilon,\delta<1$, and let $\mathcal N\subseteq \mathbb{B}_{c_r}$ be an $\eta$-net for $\mathbb{B}_{c_r}$ in the sampling seminorm $\snorm{\cdot}{\mu_{c_s}}$, where $0<\eta<\sqrt{1-\varepsilon}$. If
  \[
    m \gtrsim \varepsilon^{-2}\Lambda\left(c_r,c_r,c_s\right)\log\left(\frac{2\left|\mathcal N\right|}{\delta}\right),
  \]
  then with probability at least $1-\delta$,
  \[
    \left(\sqrt{1-\varepsilon}-\eta\right)\left\|\mathbf{h}\right\|_2
    \leq
    \left\|\mathbf{A}_\Omega \mathbf{h}\right\|_2
    \leq
    \left(\sqrt{1+\varepsilon}+\eta\right)\left\|\mathbf{h}\right\|_2,
    \qquad
    \mathbf{h}\in \mathbb{F}_{c_r}-\mathbb{F}_{c_r}.
  \]
\end{theorem}

Before proving the theorem, we record the auxiliary lemmas used in the net-extension argument.

\begin{lemma}
  \label{lem:seminorm-domination}
  For every $\mathbf{g}\in\mathbb{X}$, we have that 
  $
    \left\|\mathbf{A}_\Omega\mathbf{g}\right\|_2
    \leq
    \snorm{\mathbf{g}}{\mu_{c_s}}.
  $
\end{lemma}

\begin{proof}
  For each sampled index $I_j$,
  \[
    \frac{\left|\mathbf{P}_{I_j}\mathbf{F}\mathbf{g}\right|^2}{\mu_{c_s}\left(I_j\right)}
    \leq 
    \snorm{\mathbf g}{\mu_{c_s}}^2.
  \]
  Averaging over $j=1,\dots,m$ gives
  \[
    \left\|\mathbf{A}_\Omega\mathbf{g}\right\|_2^2
    =
    \frac{1}{m}\sum_{j=1}^m
    \frac{\left|\mathbf{P}_{I_j}\mathbf{F}\mathbf{g}\right|^2}{\mu_{c_s}\left(I_j\right)}
    \leq
    \snorm{\mathbf g}{\mu_{c_s}}^2.
  \]
  Taking square roots yields the claim.
\end{proof}

Lemma~\ref{lem:seminorm-domination} is the basic comparison that lets us pass from the sampling seminorm to the actual measurement operator. The following lemma also shows how the sampling seminorm is controlled by the Euclidean norm.

\begin{lemma}
    \label{lem:active-support-seminorm-comparison}
  For every $\mathbf g\in \mathbb X$, we have that
  \[
    \snorm{\mathbf g}{\mu_{c_s}}
    \leq
    \underline{\mu}^{-1/2}
    \left\|\mathbf g\right\|_2.
  \]
\end{lemma}


\begin{proof}
  Let $\mathbf g \in \mathbb X$. Then
  \[
    \snorm{\mathbf g}{\mu_{c_s}}^2
    =
    \max_{i \in D}
    \frac{\left|\mathbf P_i\mathbf F\mathbf g\right|^2}{\mu_{c_s}\left(i\right)}
    \leq
    \underline{\mu}^{-1}
    \max_{i \in D}
    \left|\mathbf P_i\mathbf F\mathbf g\right|^2.
  \]
  Using $\max_i a_i^2 \leq \sum_i a_i^2$ and that $\mathbf{F}$ is the discrete unitary Fourier transform, Parseval's theorem gives us
  \[
    \snorm{\mathbf g}{\mu_{c_s}}^2
    \leq
    \underline{\mu}^{-1}
    \sum_{i \in D}
    \left|\mathbf P_i\mathbf F\mathbf g\right|^2
    =
    \underline{\mu}^{-1}
    \left\|\mathbf g\right\|_2^2.
  \]
  Taking square roots gives the claim.
\end{proof}

We now record a fixed-direction concentration estimate for the weighted subsampled Fourier operator, which will later be combined with a union bound over the covering net.

\begin{lemma}
  \label{lem:fixed-direction}
  Let $\mathbf{h}\in \left(\mathbb{F}_{c_r}-\mathbb{F}_{c_r}\right)\setminus\left\{\mathbf{0}\right\}$, and let $0<\varepsilon<1$. Then
  \[
    \mathbb{P}\left(
      \left|
      \left\|\mathbf{A}_\Omega\mathbf{h}\right\|_2^2
      -
      \left\|\mathbf{h}\right\|_2^2
      \right|
      >
      \varepsilon\left\|\mathbf{h}\right\|_2^2
    \right)
    \leq
    2\exp\left(-a\,\frac{m\varepsilon^2}{\Lambda\left(c_r,c_r,c_s\right)}\right).
  \]
  for some absolute constant $a > 0$.
\end{lemma}

\begin{proof}
  Set
  \[
    X_j
    =
    \frac{\left|\mathbf{P}_{I_j}\mathbf{F}\mathbf{h}\right|^2}{\mu_{c_s}\left(I_j\right)},
    \qquad
    j=1,\dots,m.
  \]
  Since the indices $I_j$ are i.i.d. draws from $\mu_{c_s}$, the random variables $X_1,\dots,X_m$ are i.i.d., and
  \[
    \left\|\mathbf{A}_\Omega\mathbf{h}\right\|_2^2
    =
    \frac{1}{m}\sum_{j=1}^m X_j.
  \]
Since $\mu_{c_s}$ has full support,
\[
  \mathbb E[X_j]
  =
  \sum_{i\in D}
  \mu_{c_s}(i)
  \frac{|\mathbf P_i\mathbf F\mathbf h|^2}{\mu_{c_s}(i)}
  =
  \sum_{i\in D}
  |\mathbf P_i\mathbf F\mathbf h|^2
  =
  \|\mathbf h\|_2^2.
\]

  Now since $\mathbf{h}\in\mathbb{F}_{c_r}-\mathbb{F}_{c_r}$, the definition of $K(\mathbb{F}_{c_r}-\mathbb{F}_{c_r})(i)$ gives
  \[
    \frac{\left|\mathbf{P}_i\mathbf{F}\mathbf{h}\right|^2}{\left\|\mathbf{h}\right\|_2^2}
    \leq
    K(\mathbb{F}_{c_r}-\mathbb{F}_{c_r})(i),
    \qquad
    i\in D.
  \]
  Dividing by $\mu_{c_s}\left(i\right)$ yields
  \[
    \frac{\left|\mathbf{P}_i\mathbf{F}\mathbf{h}\right|^2}{\mu_{c_s}\left(i\right)}
    \leq
    \frac{K(\mathbb{F}_{c_r}-\mathbb{F}_{c_r})(i)}{\mu_{c_s}\left(i\right)}
    \left\|\mathbf{h}\right\|_2^2
    \leq
    \Lambda\left(c_r,c_r,c_s\right)\left\|\mathbf{h}\right\|_2^2.
  \]
  Hence
  \[
    0\leq X_j\leq \Lambda\left(c_r,c_r,c_s\right)\left\|\mathbf{h}\right\|_2^2.
  \]
  Let
  \[
    M=\Lambda\left(c_r,c_r,c_s\right)\left\|\mathbf{h}\right\|_2^2.
  \]
  Since $0\leq X_j\leq M$, one has $X_j^2\leq M X_j$, and therefore
  \[
    \mathbb{E}\left[X_j^2\right]
    \leq
    M\,\mathbb{E}\left[X_j\right]
    =
    \Lambda\left(c_r,c_r,c_s\right)\left\|\mathbf{h}\right\|_2^4.
  \]
  Thus
  \[
    \operatorname{Var}\left(X_j\right)
    \leq
    \Lambda\left(c_r,c_r,c_s\right)\left\|\mathbf{h}\right\|_2^4.
  \]

  Now we apply Bernstein's inequality to
  \[
    Y_j=X_j-\mathbb{E}\left[X_j\right].
  \]
  These random variables are i.i.d., mean zero, satisfy $\left|Y_j\right|\leq M$, and have variance bounded as above. Therefore
  \[
    \mathbb{P}\left(
      \left|
      \frac{1}{m}\sum_{j=1}^m Y_j
      \right|
      >
      \varepsilon\left\|\mathbf{h}\right\|_2^2
    \right)
    \leq
    2\exp\left(-a\,\frac{m\varepsilon^2}{\Lambda\left(c_r,c_r,c_s\right)}\right),
  \]
  for some absolute constant $a > 0$, which is exactly the claim.
\end{proof}

Finally, we show that control on a finite net extends to the entire normalized secant set, provided the net is taken in the sampling seminorm. That is to say, once the sampling operator $\mathbf A_\Omega$ is controlled on a net, this lemma extends the same bounds to the full class by using closeness in the sampling seminorm.

\begin{lemma}
  \label{lem:net-extension}
  Let $\mathbb{B}\subseteq \left\{\mathbf{h}\in\mathbb{X}:\left\|\mathbf{h}\right\|_2=1\right\}$ and let $\mathcal N\subseteq \mathbb{B}$ be an $\eta$-net for $\mathbb{B}$ in the seminorm $\snorm{\cdot}{\mu_{c_s}}$. Suppose that for $0<\varepsilon<1$,
  \[
    \sqrt{1-\varepsilon}
    \leq
    \left\|\mathbf{A}_\Omega\mathbf{u}\right\|_2
    \leq
    \sqrt{1+\varepsilon},
    \qquad
    \mathbf{u}\in\mathcal N.
  \]
  Then, for every $\mathbf{h}\in\mathbb{B}$,
  \[
    \left(\sqrt{1-\varepsilon}-\eta\right)
    \leq
    \left\|\mathbf{A}_\Omega\mathbf{h}\right\|_2
    \leq
    \left(\sqrt{1+\varepsilon}+\eta\right).
  \]
\end{lemma}


\begin{proof}
  Fix $\mathbf{h}\in\mathbb{B}$. Since $\mathcal N$ is an $\eta$-net for $\mathbb{B}$ in the sampling seminorm, there exists some $\mathbf{u}\in\mathcal N$ such that
  \[
    \snorm{\mathbf h - \mathbf u}{\mu_{c_s}}\leq \eta.
  \]
  By Lemma~\ref{lem:seminorm-domination},
  \[
    \left\|\mathbf{A}_\Omega\left(\mathbf{h}-\mathbf{u}\right)\right\|_2\leq \eta.
  \]
  Since $\mathbf{u}\in\mathcal N$,
  \[
    \sqrt{1-\varepsilon}
    \leq
    \left\|\mathbf{A}_\Omega\mathbf{u}\right\|_2
    \leq
    \sqrt{1+\varepsilon}.
  \]
  The triangle and reverse-triangle inequalities now yield
  \[
    \left\|\mathbf{A}_\Omega\mathbf{h}\right\|_2
    \leq
    \sqrt{1+\varepsilon}+\eta
    \qquad\text{and}\qquad
    \left\|\mathbf{A}_\Omega\mathbf{h}\right\|_2
    \geq
    \sqrt{1-\varepsilon}-\eta,
  \]
  which gives us the desired bounds.
\end{proof}

We are now equipped to prove Theorem~\ref{thm:secant-stability}, which follows by combining fixed-direction concentration on the net with the extension step above.

\begin{proof}[Proof of Theorem~\ref{thm:secant-stability}]
  For each point $\mathbf u\in\mathcal N$, define the event
  \[
    E_{\mathbf u}
    =
    \left\{
      \left|
      \left\|\mathbf A_\Omega\mathbf u\right\|_2^2
      -
      \left\|\mathbf u\right\|_2^2
      \right|
      \leq
      \varepsilon\left\|\mathbf u\right\|_2^2
    \right\}.
  \]
  Since $\mathcal N\subseteq \mathbb B_{c_r}$, for each
  $\mathbf u\in\mathcal N$ there exists
  \[
    \mathbf s_{\mathbf u}
    \in
    \left(\mathbb F_{c_r}-\mathbb F_{c_r}\right)
    \setminus\{\mathbf 0\}
  \]
  such that
  \[
    \mathbf u
    =
    \frac{\mathbf s_{\mathbf u}}
    {\left\|\mathbf s_{\mathbf u}\right\|_2}.
  \]
  Applying Lemma~\ref{lem:fixed-direction} to $\mathbf s_{\mathbf u}$ gives
      \begin{equation}\label{eq:fixed-dir-event}
      \mathbb P\left(
        \left|
        \left\|\mathbf A_\Omega \mathbf s_{\mathbf u}\right\|_2^2
        -
        \left\|\mathbf s_{\mathbf u}\right\|_2^2
        \right|
        >
        \varepsilon\left\|\mathbf s_{\mathbf u}\right\|_2^2
      \right)
      \leq
      2\exp\left(
        -a\,\frac{m\varepsilon^2}{\Lambda(c_r,c_r,c_s)}
      \right)
    \end{equation}
  for some absolute constant $a>0$. Since
  \[
    \left\|\mathbf A_\Omega \mathbf u\right\|_2^2
    =
    \frac{\left\|\mathbf A_\Omega \mathbf s_{\mathbf u}\right\|_2^2}
    {\left\|\mathbf s_{\mathbf u}\right\|_2^2},
    \qquad
    \left\|\mathbf u\right\|_2=1,
  \]
  the event in Equation~\eqref{eq:fixed-dir-event} is exactly $E_{\mathbf u}^c$. Therefore
  \[
    \mathbb P\left(E_{\mathbf u}^c\right)
    \leq
    2\exp\left(
      -a\,\frac{m\varepsilon^2}{\Lambda(c_r,c_r,c_s)}
    \right).
  \]

  Taking a union bound over $\mathcal N$ and using the lower bound on $m$,
  we find that with probability at least $1-\delta$,
  \[
    \sqrt{1-\varepsilon}
    \leq
    \left\|\mathbf A_\Omega\mathbf u\right\|_2
    \leq
    \sqrt{1+\varepsilon},
    \qquad
    \mathbf u\in\mathcal N.
  \]
  Applying Lemma~\ref{lem:net-extension} with
  $\mathbb B=\mathbb B_{c_r}$ gives
  \begin{equation}\label{eq: Aw bounded}
    \sqrt{1-\varepsilon}-\eta
    \leq
    \left\|\mathbf A_\Omega\mathbf w\right\|_2
    \leq
    \sqrt{1+\varepsilon}+\eta,
    \qquad
    \mathbf w\in\mathbb B_{c_r}.
  \end{equation}

  Now let $\mathbf h \in (\mathbb F_{c_r}-\mathbb F_{c_r}) \setminus \{ \mathbf 0 \}$ and set
  $
    \mathbf w
    =
    \mathbf h / \left\|\mathbf h\right\|_2.
  $
  Then $\mathbf w\in\mathbb B_{c_r}$, Equation~\eqref{eq: Aw bounded} gives
  \[
    \left(\sqrt{1-\varepsilon}-\eta\right)\left\|\mathbf h\right\|_2
    \leq
    \left\|\mathbf A_\Omega \mathbf h\right\|_2
    \leq
    \left(\sqrt{1+\varepsilon}+\eta\right)\left\|\mathbf h\right\|_2.
  \]
  This proves the theorem.
\end{proof}

Finally, we can prove the general Christoffel sample complexity statement in Theorem~\ref{thm:christoffel-srec}.

\begin{proof}[Proof of Theorem~\ref{thm:christoffel-srec}]
  Apply Theorem~\ref{thm:secant-stability} with $\varepsilon=\tau/2$. Then, with probability at least $1-\delta$,
  \[
    \left(\sqrt{1-\tau/2}-\eta\right)\left\|\mathbf{h}\right\|_2
    \leq
    \left\|\mathbf{A}_\Omega \mathbf{h}\right\|_2
    \leq
    \left(\sqrt{1+\tau/2}+\eta\right)\left\|\mathbf{h}\right\|_2,
    \qquad
    \mathbf{h}\in \mathbb{F}_{c_r}-\mathbb{F}_{c_r}.
  \]

  For the lower bound,
  \[
    \left\|\mathbf{A}_\Omega \mathbf{h}\right\|_2^2
    \geq
    \left(\sqrt{1-\tau/2}-\eta\right)^2\left\|\mathbf{h}\right\|_2^2
    \geq
    \left(1-\tau/2-2\eta\right)\left\|\mathbf{h}\right\|_2^2
    \geq
    \left(1-\tau\right)\left\|\mathbf{h}\right\|_2^2,
  \]
  since $\eta\leq \tau/8$.

  For the upper bound,
  \[
    \left\|\mathbf{A}_\Omega \mathbf{h}\right\|_2^2
    \leq
    \left(\sqrt{1+\tau/2}+\eta\right)^2\left\|\mathbf{h}\right\|_2^2
    =
    \left(
      1+\tau/2+2\eta\sqrt{1+\tau/2}+\eta^2
    \right)\left\|\mathbf{h}\right\|_2^2.
  \]
  Since $0<\tau<1$, we have $\sqrt{1+\tau/2}\leq \sqrt{3/2}<5/4$, and therefore
  \[
    2\eta\sqrt{1+\tau/2}+\eta^2
    \leq
    \frac{5}{2}\eta+\eta^2
    \leq
    \frac{5\tau}{16}+\frac{\tau^2}{64}
    <
    \frac{\tau}{2}.
  \]
  Hence
  \[
    \left\|\mathbf{A}_\Omega \mathbf{h}\right\|_2^2
    \leq
    \left(1+\tau\right)\left\|\mathbf{h}\right\|_2^2.
  \]
  Thus $\mathbf A_\Omega$ is empirically nondegenerate on $\mathbb F_{c_r}-\mathbb F_{c_r}$ with distortion $\tau$. Applying the lower half of this estimate to $\mathbf h=\mathbf f_1-\mathbf f_2$ for arbitrary $\mathbf f_1,\mathbf f_2\in\mathbb F_{c_r}$ yields that $\mathbf A_\Omega$ satisfies
  $
  \mathrm{S\text{-}REC}\left(\mathbb{F}_{c_r},\sqrt{1-\tau},0\right).
  $
\end{proof}


\subsection{A ReLU Specialization}\label{app:relu-specialization}

We now formalize the piecewise-linear structural assumption used in the main text and prove the corresponding sample complexity theorem. The result follows under the notion that the activation patterns of a bias-free ReLU network partition the latent space into finitely many polyhedral cones, and on each such cone the network is
linear~\cite{bora2017compressed, berk2023modeladapted, berk2022fourier, plan2025denoising, naderi2021beyond}. Motivated by this structure, we define the following piecewise-linear hypothesis:


\begin{definition}[$(N,k)$-Piecewise Linear]
  \label{def:generator-structure}
  Fix a recovery prompt $c_r$.
  We say that $G\left(\cdot,c_r\right)$ is \emph{$(N,k)$-piecewise linear} on $B_2^k(R)$ if there exist polyhedral cones $Q_1,\dots,Q_{N}\subseteq \mathbb{R}^k$ and linear maps $A_1,\dots,A_{N}:\mathbb{R}^k\to\mathbb{X}$ such that
  \[
    B_2^k(R)
    \subseteq
    \bigcup_{j=1}^{N} Q_j,
    \qquad
    G\left(\mathbf z,c_r\right)=A_j\mathbf z,
    \quad
    \mathbf z\in B_2^k(R)\cap Q_j.
  \]
\end{definition}

This structural assumption is useful because it reduces the geometry of the generator range to a finite union of low-dimensional linear subspaces. Consequently, the normalized secant set can be covered by combining Euclidean nets on at most $2k$-dimensional subspaces, leading to the following S-REC guarantee:

\begin{theorem}[Sample Complexity for Piecewise Linear Networks]
  \label{thm:piecewise-linear-srec}
  Assume that $G\left(\cdot,c_r\right)$ is $(N,k)$-piecewise linear on $B_2^k(R)$. Fix $0<\tau,\delta<1$. If
  \[
    m
    \gtrsim
    \tau^{-2}
    \Lambda\left(c_r,c_r,c_s\right)
    \left[
      \log\left(N\right)
      +
      k\log\left(
        1+
        \frac{1}
        {\tau\sqrt{\underline{\mu}}}
      \right)
      +
      \log\left(\frac{1}{\delta}\right)
    \right],
  \]
  then with probability at least $1-\delta$ over the draw of $\Omega$, $\mathbf A_\Omega$ satisfies
  $
  \mathrm{S\text{-}REC}
  \left(
    \mathbb{F}_{c_r},
    \sqrt{1-\tau},
    0
  \right).
  $
\end{theorem}

The key geometric input for Theorem~\ref{thm:piecewise-linear-srec} is the following covering estimate for the normalized secant set induced by a piecewise-linear generator.

\begin{proposition}
  \label{prop:piecewise-linear-cover-transfer}
  Assume that $G(\cdot,c_r)$ is $(N,k)$-piecewise linear on
  $B_2^k(R)$ in the sense of Definition~\ref{def:generator-structure}.. Then for every $\eta>0$, there
  exists a set $\mathcal N\subseteq\mathbb B_{c_r}$ such that $\mathcal N$ is
  an $\eta$-net for $\mathbb B_{c_r}$ in the sampling seminorm
  $\snorm{\cdot}{\mu_{c_s}}$ and
  \[
    |\mathcal N|
    \leq
    N^2
    \left(
      1+
      \frac{2}
      {\eta\sqrt{\underline{\mu}}}
    \right)^{2k}.
  \]
\end{proposition}

\begin{proof}
  Let $Q_1,\dots,Q_{N}$ and
  $A_1,\dots,A_{N}$ be as in
  Definition~\ref{def:generator-structure}. Define
  \[
    C_j=A_j(Q_j),
    \qquad
    j=1,\dots,N.
  \]
  Since each $Q_j$ is a cone and each $A_j$ is linear, each $C_j$ is a cone
  contained in a subspace of dimension at most $k$. Moreover,
  \[
    \mathbb F_{c_r}
    \subseteq
    \bigcup_{j=1}^{N} C_j.
  \]

  For each pair $(a,b)$, define
  \[
    D_{a,b}=C_a-C_b,
    \qquad
    L_{a,b}=\operatorname{span}(D_{a,b}).
  \]
  Then
  \[
    \mathbb F_{c_r}-\mathbb F_{c_r}
    \subseteq
    \bigcup_{a,b=1}^{N}D_{a,b},
  \]
  and
  \[
    \dim L_{a,b}\leq 2k.
  \]

  Now define
  \[
    E_{a,b}
    =
    \mathbb B_{c_r}\cap D_{a,b}.
  \]
  Then
  \[
    \mathbb B_{c_r}
    \subseteq
    \bigcup_{a,b=1}^{N} E_{a,b}.
  \]
  Indeed, if $\mathbf u\in\mathbb B_{c_r}$, then
  $\mathbf u=\mathbf h/\|\mathbf h\|_2$ for some nonzero
  $\mathbf h\in\mathbb F_{c_r}-\mathbb F_{c_r}$. Since
  $\mathbf h\in D_{a,b}$ for some pair $(a,b)$ and $D_{a,b}$ is a cone, also
  $\mathbf u\in D_{a,b}$. Thus $\mathbf u\in E_{a,b}$.

  For each pair $(a,b)$ with $E_{a,b}\neq\emptyset$, let
  $\mathcal N_{a,b}\subseteq E_{a,b}$ be a maximal
  $\eta\sqrt{\underline\mu}$-separated subset of $E_{a,b}$ in
  Euclidean norm. Then $\mathcal N_{a,b}$ is an
  $\eta\sqrt{\underline\mu}$-net for $E_{a,b}$ in Euclidean norm.

  Since $E_{a,b}$ is contained in the unit sphere of the subspace $L_{a,b}$
  and $\dim L_{a,b}\leq 2k$, the standard volumetric packing bound gives
  \[
    |\mathcal N_{a,b}|
    \leq
    \left(
      1+
      \frac{2}
      {\eta\sqrt{\underline\mu}}
    \right)^{2k}.
  \]

  Define
  \[
    \mathcal N
    =
    \bigcup_{a,b=1}^{N}\mathcal N_{a,b},
  \]
  omitting empty sets. Since each $\mathcal N_{a,b}\subseteq E_{a,b}$ and
  each $E_{a,b}\subseteq\mathbb B_{c_r}$, we have
  $
    \mathcal N\subseteq\mathbb B_{c_r}.
  $
  Also,
  \[
    |\mathcal N|
    \leq
    N^2
    \left(
      1+
      \frac{2}
      {\eta\sqrt{\underline\mu}}
    \right)^{2k}.
  \]

  It remains to show that $\mathcal N$ is an $\eta$-net for
  $\mathbb B_{c_r}$ in $\snorm{\cdot}{\mu_{c_s}}$. Fix
  $\mathbf u\in\mathbb B_{c_r}$. Then $\mathbf u\in E_{a,b}$ for some
  $(a,b)$, so there exists $\mathbf v\in\mathcal N_{a,b}$ such that
  \[
    \|\mathbf u-\mathbf v\|_2
    \leq
    \eta\sqrt{\underline\mu}.
  \]
  By Lemma~\ref{lem:active-support-seminorm-comparison},
  \[
    \snorm{\mathbf u-\mathbf v}{\mu_{c_s}}
    \leq
    \underline\mu^{-1/2}
    \|\mathbf u-\mathbf v\|_2
    \leq
    \eta.
  \]
  Hence $\mathcal N$ is an $\eta$-net for $\mathbb B_{c_r}$ in
  $\snorm{\cdot}{\mu_{c_s}}$.
\end{proof}

With the covering estimate in hand, the piecewise-linear theorem is obtained by feeding that cover size into the general Christoffel sampling S-REC theorem from Appendix~\ref{app:cs4ml}.

\begin{proof}[Proof of Theorem~\ref{thm:piecewise-linear-srec}]
  Set
  $
    \eta=\tau/8.
  $
  By Proposition~\ref{prop:piecewise-linear-cover-transfer}, there exists an
  $\eta$-net $\mathcal N\subseteq\mathbb B_{c_r}$ for $\mathbb B_{c_r}$ in
  the sampling seminorm such that
  \[
    |\mathcal N|
    \leq
    N^2
    \left(
      1+
      \frac{16}
      {\tau\sqrt{\underline{\mu}}}
    \right)^{2k}.
  \]
  Therefore
  \[
    \log|\mathcal N|
    \leq
    2\log N
    +
    2k
    \log\left(
      1+
      \frac{16}
      {\tau\sqrt{\underline{\mu}}}
    \right).
  \]
  Hence the stated lower bound on $m$ implies, after adjusting the absolute
  constant hidden in $\gtrsim$, that
  \[
    m
    \gtrsim
    \tau^{-2}
    \Lambda(c_r,c_r,c_s)
    \log\left(\frac{2|\mathcal N|}{\delta}\right).
  \]
  Applying Theorem~\ref{thm:christoffel-srec} gives that, with probability
  at least $1-\delta$, $\mathbf A_\Omega$ satisfies
  $
    \mathrm{S\text{-}REC}
    \left(
      \mathbb F_{c_r},
      \sqrt{1-\tau},
      0
    \right).
  $
\end{proof}

We now verify the ReLU sample complexity theorem from the main text by applying Theorem~\ref{thm:piecewise-linear-srec}, but first introduce one technical lemma from~\cite[Lemma S2.2]{berk2022fourier} (see also \cite{naderi2021beyond}, to which we refer the reader for the proof). 

\begin{lemma}[Finite Cone Decomposition for ReLU Generators]
  \label{lem:relu-cone-decomposition}
  Assume that, for the fixed prompt $c_r$, the map $G\left(\cdot,c_r\right)$ is a bias-free depth-$d$ feedforward ReLU network in the sense of Definition~\ref{def:bias-free-relu} with layer widths
  $
    k = k_0 \leq k_1,\dots,k_{d-1},
    k_d = n
  $,
  with $d \geq 2$, and define
  \[
    \overline{k}_{c_r}
    =
    \left(
      \prod_{\ell=1}^{d-1} k_\ell
    \right)^{1/(d-1)}.
  \]
  Then $G\left(\cdot,c_r\right)$ is $(N,k)$-piecewise linear on $B_2^k(R)$ in the sense of Definition~\ref{def:generator-structure}, with
  \[
      \log N
      \leq
      k(d-1)\log\left(
        \frac{2e\,\overline{k}_{c_r}}{k}
      \right).
    \]
\end{lemma}

Once the ReLU network has been rewritten as a piecewise-linear model with an explicit region count we are ready to prove Theorem~\ref{thm:relu-srec}.

\begin{proof}[Proof of Theorem~\ref{thm:relu-srec}]
  Lemma~\ref{lem:relu-cone-decomposition} verifies the hypotheses of Theorem~\ref{thm:piecewise-linear-srec} and supplies the stated bound on $\log N$. Applying Theorem~\ref{thm:piecewise-linear-srec} gives the claim.
\end{proof}

\subsection{A Lipschitz specialization}\label{app:lipschitz-specialization}

We now prove the Lipschitz specialization stated in the main text. Throughout this subsection, assume that $G\left(\cdot,c_r\right)$ is $L$-Lipschitz on $B_2^k\left(R\right)$. That is to say 
  \[
    \left\|
    G\left(\mathbf z_1,c_r\right)
    -
    G\left(\mathbf z_2,c_r\right)
    \right\|_2
    \leq
    L\left\|\mathbf z_1-\mathbf z_2\right\|_2,
    \qquad
    \mathbf z_1,\mathbf z_2\in B_2^k(R).
  \] 
The Euclidean covering argument below is the standard image-of-a-net construction with truncation on $B_2^k\left(R\right)$, similar to that in \cite[Lemma 4.1 and its proof]{bora2017compressed}. The only additional step is to convert the Euclidean cover into one measured in the sampling seminorm. That is to say we first build a Euclidean cover in latent space, and then transfer that cover to the sampling seminorm.

\begin{proposition}
  \label{prop:lipschitz-cover-transfer}
  Assume that $G(\cdot,c_r)$ is $L$-Lipschitz on $B_2^k(R)$. Let $0<\eta\leq 1$ and $\xi>0$. Then
  there exists a set $\mathcal N\subseteq\mathbb B_{c_r}$ such that
  $\mathcal N$ is an $\eta$-net for $\mathbb B_{c_r,\xi}$ in the sampling
  seminorm $\snorm{\cdot}{\mu_{c_s}}$ and
  \[
    |\mathcal N|
    \leq
    \left(
      1+
      \frac{8LR}
      {\xi\eta\sqrt{\underline{\mu}}}
    \right)^{2k}.
  \]
\end{proposition}

\begin{proof}
  Set
  \[
    s=\frac{\xi\eta\sqrt{\underline{\mu}}}{2}.
  \]
  Since $0<\eta\leq 1$ and $\underline{\mu}\leq 1$, we have
  $s\leq \xi/2$.

  We first build an explicit Euclidean $s$-net for
  $\mathbb F_{c_r}-\mathbb F_{c_r}$ with centers inside
  $\mathbb F_{c_r}-\mathbb F_{c_r}$. Let $\mathcal Z$ be an
  $s/(2L)$-net for $B_2^k(R)$ in Euclidean norm. By the standard
  volumetric bound,
  \[
    |\mathcal Z|
    \leq
    \left(
      1+
      \frac{4LR}{s}
    \right)^k.
  \]
  Define
  \[
    \mathcal H
    =
    \left\{
      G(\mathbf z_1,c_r)-G(\mathbf z_2,c_r)
      :
      \mathbf z_1,\mathbf z_2\in\mathcal Z
    \right\}.
  \]
  Then
  \[
    |\mathcal H|
    \leq
    |\mathcal Z|^2
    \leq
    \left(
      1+
      \frac{4LR}{s}
    \right)^{2k}.
  \]
  Moreover, $\mathcal H$ is an $s$-net for
  $\mathbb F_{c_r}-\mathbb F_{c_r}$ in Euclidean norm. Indeed, for any
  \[
    \mathbf h
    =
    G(\mathbf z_1,c_r)-G(\mathbf z_2,c_r)
    \in
    \mathbb F_{c_r}-\mathbb F_{c_r},
  \]
  choose $\tilde{\mathbf z}_1,\tilde{\mathbf z}_2\in\mathcal Z$ such that
  \[
    \|\mathbf z_j-\tilde{\mathbf z}_j\|_2
    \leq
    \frac{s}{2L},
    \qquad
    j=1,2.
  \]
  Then, by the Lipschitz property,
  \[
    \left\|
      \mathbf h
      -
      \left(
        G(\tilde{\mathbf z}_1,c_r)
        -
        G(\tilde{\mathbf z}_2,c_r)
      \right)
    \right\|_2
    \leq
    L\|\mathbf z_1-\tilde{\mathbf z}_1\|_2
    +
    L\|\mathbf z_2-\tilde{\mathbf z}_2\|_2
    \leq
    s.
  \]

  Now define
  \[
    \mathcal N
    =
    \left\{
      \frac{\mathbf g}{\left\|\mathbf g\right\|_2}
      :
      \mathbf g\in\mathcal H,\ \mathbf g\neq \mathbf 0
    \right\}.
  \]
  Since every nonzero $\mathbf g\in\mathcal H$ lies in
  $\mathbb F_{c_r}-\mathbb F_{c_r}$, we have
  \[
    \mathcal N\subseteq\mathbb B_{c_r}.
  \]

  We claim that $\mathcal N$ covers $\mathbb B_{c_r,\xi}$ in
  $\snorm{\cdot}{\mu_{c_s}}$ with radius $\eta$. Fix
  $\mathbf u\in\mathbb B_{c_r,\xi}$. Then
  \[
    \mathbf u=\frac{\mathbf h}{\|\mathbf h\|_2}
  \]
  for some $\mathbf h\in\mathbb F_{c_r}-\mathbb F_{c_r}$ with
  $\left\|\mathbf h \right\|_2\geq \xi$. Since $\mathcal H$ is an $s$-net for
  $\mathbb F_{c_r}-\mathbb F_{c_r}$, choose $\mathbf g\in\mathcal H$ such
  that
  \[
    \left\|\mathbf h-\mathbf g\right\|_2\leq s.
  \]
  Since $s\leq \xi/2$ and $\left\|\mathbf h\right\|_2\geq\xi$, we have
  \[
    \left\|\mathbf g\right\|_2
    \geq
    \left\|\mathbf h\right\|_2-\left\|\mathbf h-\mathbf g\right\|_2
    \geq
    \xi-s
    \geq
    \xi/2
    >
    0.
  \]
  Hence $\mathbf v=\mathbf g/\left\|\mathbf g\right\|_2$ belongs to $\mathcal N$.
  
We next estimate the Euclidean distance between the normalized vectors
\[
 \left\|\mathbf u-\mathbf v\right\|_2 = 
  \left\|
    \frac{\mathbf h}{\left\|\mathbf h\right\|_2}
    -
    \frac{\mathbf g}{\left\|\mathbf g\right\|_2}
  \right\|_2
  \leq
  \left\|
    \frac{\mathbf h-\mathbf g}{\left\|\mathbf h\right\|_2}
  \right\|_2
  +
  \left\|
    \frac{\mathbf g}{\left\|\mathbf h\right\|_2}
    -
    \frac{\mathbf g}{\left\|\mathbf g\right\|_2}
  \right\|_2.
\]
The first term is $\left\|\mathbf h-\mathbf g\right\|_2/\left\|\mathbf h\right\|_2$.
For the second term,
\[
  \left\|
    \frac{\mathbf g}{\left\|\mathbf h\right\|_2}
    -
    \frac{\mathbf g}{\left\|\mathbf g\right\|_2}
  \right\|_2
  =
  \left\|\mathbf g\right\|_2
  \left|
    \frac{1}{\left\|\mathbf h\right\|_2}
    -
    \frac{1}{\left\|\mathbf g\right\|_2}
  \right|
  =
  \frac{\left|\left\|\mathbf g\right\|_2-\left\|\mathbf h\right\|_2\right|}
       {\left\|\mathbf h\right\|_2}
  \leq
  \frac{\left\|\mathbf h-\mathbf g\right\|_2}{\left\|\mathbf h\right\|_2}.
\]
Therefore
\[
  \left\|\mathbf u-\mathbf v\right\|_2
  \leq
  \frac{2\left\|\mathbf h-\mathbf g\right\|_2}{\left\|\mathbf h\right\|_2}
  \leq
  \frac{2s}{\xi}
  =
  \eta\sqrt{\underline{\mu}}.
\]

By Lemma~\ref{lem:active-support-seminorm-comparison},
\[
  \snorm{\mathbf u-\mathbf v}{\mu_{c_s}}
  \leq
  \underline{\mu}^{-1/2}
  \left\|\mathbf u-\mathbf v\right\|_2
  \leq
  \eta.
\]
Thus $\mathcal N$ is an $\eta$-net for $\mathbb B_{c_r,\xi}$ in
$\snorm{\cdot}{\mu_{c_s}}$. Finally, we estimate
\[
  \left|\mathcal N\right|
  \leq
  \left|\mathcal H\right|
  \leq
  \left(
    1+
    \frac{4LR}{s}
  \right)^{2k}
  =
  \left(
    1+
    \frac{8LR}
         {\xi\eta\sqrt{\underline{\mu}}}
  \right)^{2k},
\]
as desired.
\end{proof}

 With this argument, we are now ready to prove Theorem~\ref{thm:lipschitz-srec}.

\begin{proof}[Proof of Theorem~\ref{thm:lipschitz-srec}]
  Set
  $
    \eta=\tau/8.
  $
  By Proposition~\ref{prop:lipschitz-cover-transfer}, there exists a set
  $\mathcal N\subseteq\mathbb B_{c_r}$ that is an $\eta$-net for
  $\mathbb B_{c_r,\xi}$ in the sampling seminorm and satisfies
  \[
    \left|\mathcal N\right|
    \leq
    \left(
      1+
      \frac{64LR}
           {\xi\tau\sqrt{\underline{\mu}}}
    \right)^{2k}.
  \]
  Therefore the stated lower bound on $m$ implies, after adjusting the
  absolute constant hidden in $\gtrsim$, that
  \[
    m
    \gtrsim
    \left(\frac{\tau}{2}\right)^{-2}
    \Lambda\left(c_r,c_r,c_s\right)
    \log\left(
      \frac{2\left|\mathcal N\right|}{\delta}
    \right).
  \]

  Since $\mathcal N\subseteq\mathbb B_{c_r}$, every
  $\mathbf u\in\mathcal N$ is a normalized nonzero secant. Applying
  Lemma~\ref{lem:fixed-direction} with $\varepsilon=\tau/2$ and taking a
  union bound over $\mathcal N$, we obtain that, with probability at least
  $1-\delta$,
  \[
    \sqrt{1-\tau/2}
    \leq
    \left\|\mathbf A_\Omega\mathbf u\right\|_2
    \leq
    \sqrt{1+\tau/2},
    \qquad
    \mathbf u\in\mathcal N.
  \]

  We now extend the lower bound from $\mathcal N$ to
  $\mathbb B_{c_r,\xi}$. Fix $\mathbf v\in\mathbb B_{c_r,\xi}$. Since
  $\mathcal N$ is an $\eta$-net for $\mathbb B_{c_r,\xi}$ in
  $\snorm{\cdot}{\mu_{c_s}}$, there exists $\mathbf u\in\mathcal N$ such that
  \[
    \snorm{\mathbf v-\mathbf u}{\mu_{c_s}}\leq \eta.
  \]
  By Lemma~\ref{lem:seminorm-domination},
  \[
    \left\|\mathbf A_\Omega\left(\mathbf v-\mathbf u\right)\right\|_2
    \leq
    \eta.
  \]
  Therefore,
  \[
    \left\|\mathbf A_\Omega\mathbf v\right\|_2
    \geq
    \left\|\mathbf A_\Omega\mathbf u\right\|_2
    -
    \left\|\mathbf A_\Omega\left(\mathbf v-\mathbf u\right)\right\|_2
    \geq
    \sqrt{1-\tau/2}-\eta.
  \]
  Since $\eta=\tau/8$ and
  \[
    \sqrt{1-\tau/2}-\sqrt{1-\tau}
    =
    \frac{\tau/2}
         {\sqrt{1-\tau/2}+\sqrt{1-\tau}}
    \geq
    \frac{\tau}{4},
  \]
  we have
  \[
    \sqrt{1-\tau/2}-\eta
    \geq
    \sqrt{1-\tau}
    +
    \frac{\tau}{4}
    -
    \frac{\tau}{8}
    \geq
    \sqrt{1-\tau}.
  \]
  Hence
  \begin{equation}
  \left\|\mathbf A_\Omega\mathbf v\right\|_2
  \geq
  \sqrt{1-\tau},
  \qquad
  \mathbf v\in\mathbb B_{c_r,\xi}.
  \label{eq:srec-lower-bound-on-normalized-secants}
\end{equation}

  Now let
  \[
    \mathbf h\in\mathbb F_{c_r}-\mathbb F_{c_r}.
  \]
  If $\left\|\mathbf h\right\|_2\geq\xi$, then
  $\mathbf v=\mathbf h/\left\|\mathbf h\right\|_2$ belongs to
  $\mathbb B_{c_r,\xi}$, and and Equation~\eqref{eq:srec-lower-bound-on-normalized-secants} gives
  \[
    \left\|\mathbf A_\Omega\mathbf h\right\|_2
    \geq
    \sqrt{1-\tau}\left\|\mathbf h\right\|_2.
  \]
  If instead $\left\|\mathbf h\right\|_2<\xi$, then
  \[
    \sqrt{1-\tau}\left\|\mathbf h\right\|_2
    -
    \sqrt{1-\tau}\,\xi
    \leq
    0,
  \]
  and therefore
  \[
    \left\|\mathbf A_\Omega\mathbf h\right\|_2
    \geq
    0
    \geq
    \sqrt{1-\tau}\left\|\mathbf h\right\|_2
    -
    \sqrt{1-\tau}\,\xi.
  \]
  Thus in either case,
  \[
    \left\|\mathbf A_\Omega\mathbf h\right\|_2
    \geq
    \sqrt{1-\tau}\left\|\mathbf h\right\|_2
    -
    \sqrt{1-\tau}\,\xi,
    \qquad
    \mathbf h\in\mathbb F_{c_r}-\mathbb F_{c_r}.
  \]
  Hence $\mathbf A_\Omega$ satisfies
  $
    \mathrm{S\text{-}REC}
    \left(
      \mathbb F_{c_r},
      \sqrt{1-\tau},
      \sqrt{1-\tau}\,\xi
    \right).
  $
\end{proof}

\subsection{Recovery Guarantees}\label{app:agnostic-recovery}
We now prove the deterministic recovery guarantees. Once the S-REC holds, reconstruction error is controlled by the best model approximation, the measurement residual of that approximation, the noise, and the optimization error. 

\begin{theorem}[Agnostic Recovery Under S-REC]
  \label{thm:agnostic-recovery-srec}
  Suppose that $\mathbf y=\mathbf A_\Omega\mathbf f^*+\mathbf e$. Let $c_r\in\mathcal C$ be the recovery prompt and suppose that
  $\widehat{\mathbf z}$ is an  $\omega$-minimizer in the sense of Equation~\eqref{eq: omega minimizer} with estimator $\widehat{\mathbf f}=G(\widehat{\mathbf z},c_r)$.
  If $\mathbf A_\Omega$ satisfies
  $\mathrm{S\text{-}REC}(\mathbb F_{c_r},\gamma,q)$
  for some $\gamma>0$ and $q\geq 0$, then
  \[
    \left\|\widehat{\mathbf f}-\mathbf f^*\right\|_2
    \leq
    \inf_{\mathbf f\in\mathbb F_{c_r}}
    \left[
      \left\|\mathbf f^*-\mathbf f\right\|_2
      +
      \frac{2}{\gamma}
      \left\|\mathbf A_\Omega\left(\mathbf f^*-\mathbf f\right)\right\|_2
    \right]
    +
    \frac{2\left\|\mathbf e\right\|_2+\omega+q}{\gamma}.
  \]
\end{theorem}


\begin{proof}[Proof]
  Let $\bar{\mathbf f}\in\mathbb{F}_{c_r}$ be arbitrary. Since $\widehat{\mathbf f}$ is an $\omega$-approximate weighted residual minimizer,
  \[
    \left\|
    \mathbf y
    -
    \mathbf A_\Omega\widehat{\mathbf f}
    \right\|_2
    \le
    \left\|
    \mathbf y
    -
    \mathbf A_\Omega\bar{\mathbf f}
    \right\|_2
    +
    \omega.
  \]
  Using
  \[
    \mathbf y
    =
    \mathbf A_\Omega\mathbf f^*+\mathbf e,
  \]
  this becomes
  \[
    \left\|
    \mathbf A_\Omega\left(\widehat{\mathbf f}-\mathbf f^*\right)
    -
    \mathbf e
    \right\|_2
    \le
    \left\|
    \mathbf A_\Omega\left(\bar{\mathbf f}-\mathbf f^*\right)
    -
    \mathbf e
    \right\|_2
    +
    \omega.
  \]
  Hence, by the triangle inequality,
  \[
    \left\|
    \mathbf A_\Omega\left(\widehat{\mathbf f}-\bar{\mathbf f}\right)
    \right\|_2
    \le
    2
    \left\|
    \mathbf A_\Omega\left(\bar{\mathbf f}-\mathbf f^*\right)
    \right\|_2
    +
    2\left\|\mathbf e\right\|_2
    +
    \omega.
  \]
  Since $\widehat{\mathbf f},\bar{\mathbf f}\in\mathbb{F}_{c_r}$, the S-REC assumption gives
  \[
    \gamma\left\|\widehat{\mathbf f}-\bar{\mathbf f}\right\|_2
    \le
    \left\|
    \mathbf A_\Omega\left(\widehat{\mathbf f}-\bar{\mathbf f}\right)
    \right\|_2
    +
    q.
  \]
  Therefore
  \[
    \left\|\widehat{\mathbf f}-\bar{\mathbf f}\right\|_2
    \le
    \frac{
      2
      \left\|
      \mathbf A_\Omega\left(\bar{\mathbf f}-\mathbf f^*\right)
      \right\|_2
      +
      2\left\|\mathbf e\right\|_2
      +
      \omega
      +
      q
    }{\gamma}.
  \]
  By the triangle inequality again,
  \[
    \left\|\widehat{\mathbf f}-\mathbf f^*\right\|_2
    \le
    \left\|\widehat{\mathbf f}-\bar{\mathbf f}\right\|_2
    +
    \left\|\bar{\mathbf f}-\mathbf f^*\right\|_2,
  \]
  so
  \[
    \left\|\widehat{\mathbf f}-\mathbf f^*\right\|_2
    \le
    \left\|\bar{\mathbf f}-\mathbf f^*\right\|_2
    +
    \frac{
      2
      \left\|
      \mathbf A_\Omega\left(\bar{\mathbf f}-\mathbf f^*\right)
      \right\|_2
      +
      2\left\|\mathbf e\right\|_2
      +
      \omega
      +
      q
    }{\gamma}.
  \]
  Taking the infimum over $\bar{\mathbf f}\in\mathbb{F}_{c_r}$ yields the claim.
\end{proof}

We now relate this result to the best achievable measurement-domain mismatch over $\mathbb{F}_{c_r}$ when the signal we recover is in the range of the generator.

\begin{proposition}
  \label{prop:cross-class-meas-sig}
  Let $\mathbf f^*\in \mathbb F_{c_*}$. Then, for every
  $\mathbf f\in\mathbb F_{c_r}$,
  \[
    \left\|
    \mathbf A_\Omega\left(\mathbf f^*-\mathbf f\right)
    \right\|_2
    \leq
    \sqrt{\Lambda(c_*,c_r,c_s)}
    \left\|
    \mathbf f^*-\mathbf f
    \right\|_2.
  \]
  Consequently,
  \[
    \inf_{\mathbf f\in\mathbb F_{c_r}}
    \left\|
    \mathbf A_\Omega\left(\mathbf f^*-\mathbf f\right)
    \right\|_2
    \leq
    \sqrt{\Lambda(c_*,c_r,c_s)}
    \inf_{\mathbf f\in\mathbb F_{c_r}}
    \left\|
    \mathbf f^*-\mathbf f
    \right\|_2.
  \]
\end{proposition}

\begin{proof}
  Fix $\mathbf f\in\mathbb F_{c_r}$ and set
  $
    \mathbf h=\mathbf f^*-\mathbf f.
  $
  If $\mathbf h=\mathbf 0$, the claim is immediate. Otherwise,
  since $\mathbf f^*\in\mathbb F_{c_*}$ and $\mathbf f\in\mathbb F_{c_r}$,
  we have
  $
    \mathbf h\in \mathbb F_{c_*}-\mathbb F_{c_r}.
  $
  By the definition of the cross-class Christoffel function,
  \[
    \frac{\left|\mathbf P_i\mathbf F\mathbf h\right|^2}
    {\left\|\mathbf h\right\|_2^2}
    \leq
    K\left(\mathbb F_{c_*}-\mathbb F_{c_r}\right)(i),
    \qquad i\in D.
  \]
  Therefore,
  \[
    \frac{\left|\mathbf P_i\mathbf F\mathbf h\right|^2}
    {\mu_{c_s}(i)}
    \leq
    \Lambda(c_*,c_r,c_s)\left\|\mathbf h\right\|_2^2,
    \qquad i\in D.
  \]
  Hence, for each sampled index $I_j$,
  \[
    \frac{\left|\mathbf P_{I_j}\mathbf F\mathbf h\right|^2}
    {\mu_{c_s}(I_j)}
    \leq
    \Lambda(c_*,c_r,c_s)\left\|\mathbf h\right\|_2^2.
  \]
  Averaging gives
  \[
    \left\|\mathbf A_\Omega\mathbf h\right\|_2^2
    =
    \frac{1}{m}
    \sum_{j=1}^m
    \frac{\left|\mathbf P_{I_j}\mathbf F\mathbf h\right|^2}
    {\mu_{c_s}(I_j)}
    \leq
    \Lambda(c_*,c_r,c_s)\left\|\mathbf h\right\|_2^2.
  \]
  Taking square roots proves the pointwise claim. Taking the infimum over
  $\mathbf f\in\mathbb F_{c_r}$ gives the final inequality.
\end{proof}

With this, we can now prove the recovery guarantee in Theorem~\ref{thm:prompt-mismatched-recovery}.
\begin{proof}[Proof of Theorem~\ref{thm:prompt-mismatched-recovery}]
  Combining Proposition~\ref{prop:cross-class-meas-sig} with Theorem~\ref{thm:agnostic-recovery-srec} gives us the desired result.
\end{proof}

\newpage\section{Recovered Image Comparisons}\label{app: recovered img comp}
Here we display panels of recovered images in each signal recovery task: in-range prompt-mismatched image recovery with  $c_* = \texttt{sunset over a sandy coast}$ (Section~\ref{subsec: in-range diff prompt}), out-of-range image recovery (Section~\ref{subsec: out of range}), and in-range prompt-matched image recovery with  $c_* = \texttt{"sunset beach"}$ (Appendix~\ref{app: in range prompt match exp}). All displayed images represent the best reconstruction at a sampling ratio of $0.00125$: for each fixed sampling distribution, displayed sampling ratio, and reconstruction prompt, we rank all available repeats by PSNR, breaking ties by SSIM, and display the highest-ranked reconstruction. Here, PPMAE denotes the per-pixel mean absolute error. 




%%%%%%%%%%%%%%%% In-Range Prompt-Mismatched %%%%%%%%%%%%%%
\begin{figure}[htp]
    \centering
    \includegraphics[width=1.0\linewidth]{figures/sunset_recovery_final/recovery_image_grid_k0.pdf}
    \includegraphics[width=1.0\linewidth]{figures/sunset_recovery_final/recovery_image_grid_k1_daytime_beach.pdf}
    \includegraphics[width=1.0\linewidth]{figures/sunset_recovery_final/recovery_image_grid_k2_sunset_beach.pdf}
    \includegraphics[width=1.0\linewidth]{figures/sunset_recovery_final/recovery_image_grid_k4_cat.pdf}
    \caption{Best recovered images in the in-range prompt-mismatched experiment (Section~\ref{subsec: in-range diff prompt}). Ground truth image generated with prompt $c_* = \texttt{"sunset over a sandy coast"}$.}
    \label{fig: in-range prompt-mismatch panel}
\end{figure}



%%%%%%%%%%%%%%%% Out-of-Range %%%%%%%%%%%%%%
\begin{figure}[htp]
    \centering
    \includegraphics[width=1.0\linewidth]{figures/sunset_ood_recovery_final/recovery_image_grid_k0.pdf}
    \includegraphics[width=1.0\linewidth]{figures/sunset_ood_recovery_final/recovery_image_grid_k1_daytime_beach.pdf}
    \includegraphics[width=1.0\linewidth]{figures/sunset_ood_recovery_final/recovery_image_grid_k2_sunset_beach.pdf}
    \includegraphics[width=1.0\linewidth]{figures/sunset_ood_recovery_final/recovery_image_grid_k4_cat.pdf}
    \caption{Best recovered images in the out-of-range experiment (Section~\ref{subsec: out of range}).}
    \label{fig: out of range panel}
\end{figure}






%%%%%%%%%%%%%%%% In-Range Prompt-Mismatched %%%%%%%%%%%%%%
\begin{figure}
    \centering
    \includegraphics[width=0.95\linewidth]{figures/sunset_recovery_full/recovery_image_grid_k0.pdf}
    \includegraphics[width=0.95\linewidth]{figures/sunset_recovery_full/recovery_image_grid_k1_daytime_beach.pdf}
    \includegraphics[width=0.93\linewidth]{figures/sunset_recovery_full/recovery_image_grid_k2_sunset_beach.pdf}
    \includegraphics[width=0.93\linewidth]{figures/sunset_recovery_full/recovery_image_grid_k4_cat.pdf}
    \caption{Best recovered images in the in-range prompt-matched experiment (Appendix~\ref{app: in range prompt match exp}). Ground truth image generated with prompt $c_* = \texttt{"sunset beach"}$.}
    \label{fig: in-range prompt-match panel}
\end{figure}





%%%%%%%%%%%%%%%%%%%%%%%%%%%%%%%%%%%%%%%%%%%%%%%%%%%%%%%%%%%%

\end{document}
